\section{AMC 12 12B 2023}
        \begin{enumerate}
        	\item (\textit{AMC 12 12B 2023}) Mrs. Jones is pouring orange juice into four identical glasses for her four sons. She fills the first three glasses completely but runs out of juice when the fourth glass is only $\frac{1}{3}$ full. What fraction of a glass must Mrs. Jones pour from each of the first three glasses into the fourth glass so that all four glasses will have the same amount of juice?


 $\textbf{(A) }\frac{1}{12}\qquad\textbf{(B) }\frac{1}{4}\qquad\textbf{(C) }\frac{1}{6}\qquad\textbf{(D) }\frac{1}{8}\qquad\textbf{(E) }\frac{2}{9}$

 \textbf{Answer: }C



\textbf{Solution 1}

The first three glasses each have a full glass.  Let's assume that each glass has "1 unit" of juice. It won't matter exactly how much juice everyone has because we're dealing with ratios, and that wouldn't affect our answer. The fourth glass has a glass that is one third. So the total amount of juice will be $1+1+1+\dfrac{1}{3} = \dfrac{10}{3}$. If we divide the total amount of juice by 4, we get $\dfrac{5}{6}$, which should be the amount of juice in each glass. This means that each of the first three glasses will have to contribute $1 - \dfrac{5}{6} = \boxed{\textbf{(C) }\dfrac16}$ to the fourth glass.

 



\textbf{Solution 2 (unnecessary numerical values)}

Given that the first three glasses are full and the fourth is only $\frac{1}{3}$ full, let's represent their contents with a common denominator, which we'll set as 6. This makes the first three glasses $\dfrac{6}{6}$ full, and the fourth glass $\frac{2}{6}$ full.

 To equalize the amounts, Mrs. Jones needs to pour juice from the first three glasses into the fourth. Pouring $\frac{1}{6}$ from each of the first three glasses will make them all $\dfrac{5}{6}$ full. Thus, all four glasses will have the same amount of juice. Therefore, the answer is $\boxed{\textbf{(C) }\dfrac16}.$




\textbf{Solution 3 (simple algebra)}

We let $x$ denote how much juice we take from each of the first $3$ children and give to the $4$th child. 

 We can write the following equation: $1-x=\dfrac13+3x$, since each value represents how much juice each child (equally) has in the end. (Each of the first three children now have $1-x$ juice, and the fourth child has $3x$ more juice on top of their initial $\dfrac13$.) 

 Solving, we see that $x=\boxed{\textbf{(C) }\dfrac16}.$



	\item (\textit{AMC 12 12B 2023}) Carlos went to a sports store to buy running shoes. Running shoes were on sale, with prices reduced by $20\%$on every pair of shoes. Carlos also knew that he had to pay a $7.5\%$ sales tax on the discounted price. He had $43$ dollars. What is the original (before discount) price of the most expensive shoes he could afford to buy?


 $\textbf{(A) }46\qquad\textbf{(B) }50\qquad\textbf{(C) }48\qquad\textbf{(D) }47\qquad\textbf{(E) }49$

 \textbf{Answer: }B



\textbf{Solution 1}

Let the price originally be \( x \). Then, after a \(20\) percent discount, the price is now \( x - \frac{1}{5}x = \frac{4}{5}x \). 

 From the discounted price \( \frac{4}{5}x \), we now take \( \frac{7.5}{100} \) of \( \frac{4}{5}x \) and add it to \( \frac{4}{5}x \), giving us \( \frac{4}{5}x + \left(\frac{7.5}{100}\right)\left(\frac{4}{5}x\right) = \frac{4}{5}x + \frac{30}{500}x = \frac{4}{5}x + \frac{3}{50}x = \frac{43}{50}x \). 

 Now we write the inequality \( \frac{43}{50}x \leq \$43 \) and multiply by \( \frac{50}{43} \) on both sides to get \( x \leq \$50 \). 

 We want the greatest original price, which would be \($50\) or option choice \(B\).

 



\textbf{Solution 2 (easy)}

We can create the equation:
 \[0.8x \cdot 1.075 = 43\]using the information given. This is because x, the original price, got reduced by 20\%, or multiplied by 0.8, and it also got multiplied by 1.075 on the discounted price. Solving that equation, we get
 \[\frac{4}{5} \cdot x \cdot \frac{43}{40} = 43\]
 \[\frac{4}{5} \cdot x \cdot \frac{1}{40} = 1\]
 \[\frac{1}{5} \cdot x \cdot \frac{1}{10} = 1\]
 \[x  = \boxed{50}\]




\textbf{Solution 3 (One-Step Equation)}

The discounted shoe is $20\%$ off the original price. So that means $1 - 0.2 = 0.8$. There is also a $7.5\%$ sales tax charge, so $0.8 * 1.075 = 0.86$. Now we can set up the equation $0.86x = 43$, and solving that we get $x=\boxed{\textbf{(B) }50}$ ~ kabbybear

 



\textbf{Solution 4}

Let the original price be $x$ dollars. After the discount, the price becomes $80\%\cdot x$ dollars. After tax, the price becomes $80\% \times (1+7.5\%) = 86\% \cdot x$ dollars.So, $43=86\% \cdot x$, $x=\boxed{\textbf{(B) }\$50}.$




\textbf{Solution 5}

We can assign a variable $c$ to represent the original cost of the shoes. Next, we set up the equation $80\%\cdot107.5\%\cdot c=43$. We can solve this equation for $c$ and get $\boxed{\textbf{(B) }\$50}$.

 



\textbf{Solution 6 (Intuition and Guessing)}

We know the discount price will be 5/4, and 0.075 is equal to 3/40. So we look at answer choice $\textbf{(B) }$, see that the discount price will be 40, and with sales tax applied it will be 43, so the answer choice is $\boxed{\textbf{(B) }\$50}$.

 


	\item (\textit{AMC 12 12B 2023}) A $3-4-5$ right triangle is inscribed in circle $A$, and a $5-12-13$ right triangle is inscribed in circle $B$. What is the ratio of the area of circle $A$ to the area of circle $B$?


 $\textbf{(A)}~\frac{9}{25}\qquad\textbf{(B)}~\frac{1}{9}\qquad\textbf{(C)}~\frac{1}{5}\qquad\textbf{(D)}~\frac{25}{169}\qquad\textbf{(E)}~\frac{4}{25}$

 \textbf{Answer: }D



\textbf{Solution 1}

Because the triangles are right triangles, we know the hypotenuses are diameters of circles $A$ and $B$. Thus, their radii are 2.5 and 6.5 (respectively). Square the two numbers and multiply $\pi$ to get $6.25\pi$ and $42.25\pi$ as the areas of the circles. Multiply 4 on both numbers to get $25\pi$ and $169\pi$. Cancel out the $\pi$, and lastly, divide, to get your answer $=\boxed{\textbf{(D) }\frac{25}{169}}.$




\textbf{Solution 1.5 (Way FASTER! Use in the contest)}

With right triangles inscribed in circles, the hypotenuse must be the diameter. Therefore, the ratio of radii is $\frac{5}{13}$  which means the  area ratio is just $\frac{5^2}{13^2}$  so,  $\boxed{\textbf{(D) }\frac{25}{169}}.$

 -Mismatchedcubing/Andrew\_

 \textbf{Note for Solution 1} The reason why the right triangle has its hypotenuse on the diameter is because of its properties with the circle and the central angle theorem. The central angle theorem states that an angle inscribed on a circle is just arc length in degrees. The diameter of a circle splits the circle into two halves with arc of \( 180^\circ \), and therefore the angle inscribed in that arc is just \( 180/2 = 90 \) degrees. The longest side of a triangle is opposite to the longest angle, and therefore the length of the hypotenuse is just the diameter of the circle.

 



\textbf{Solution 2}

Since the arc angle of the diameter of a circle is $90$ degrees, the hypotenuse of each these two triangles is respectively the diameter of circles $A$ and $B$. 

 Therefore the ratio of the areas equals the radius of circle $A$ squared : the radius of circle $B$ squared$=$ $0.5\times$ the diameter of circle $A$, squared : $0.5\times$ the diameter of circle $B$, squared$=$ the diameter of circle $A$, squared: the diameter of circle $B$, squared $=\boxed{\textbf{(D) }\frac{25}{169}}.$




\textbf{Solution 3}

The ratio of areas of circles is the same as the ratios of the diameters squared (since they are similar figures). Since this is a right triangle the hypotenuse of each triangle will be the diameter of the circle. This yields the expression $\frac{5^2}{13^2} =\boxed{\textbf{(D) }\frac{25}{169}}.$



	\item (\textit{AMC 12 12B 2023}) Jackson's paintbrush makes a narrow strip with a width of $6.5$ millimeters. Jackson has enough paint to make a strip $25$ meters long. How many square centimeters of paper could Jackson cover with paint?


 $\textbf{(A) }162,500\qquad\textbf{(B) }162.5\qquad\textbf{(C) }1,625\qquad\textbf{(D) }1,625,000\qquad\textbf{(E) }16,250$

 \textbf{Answer: }C



\textbf{Solution 1}

$6.5$ millimeters is equal to $0.65$ centimeters. $25$ meters is $2500$ centimeters. The answer is $0.65 \times 2500$, so the answer is $\boxed{\textbf{(C) 1,625}}$. 

 



\textbf{Solution 2 (Scientific Form)}

$6.5$ millimeters can be represented as $65 \times 10^{-2}$ centimeters. $25$ meters is $25 \times 10^{2}$ centimeters. Multiplying out these results in $(65 \times 10^{-2}) \times (25 \times 10^{2})$, which is $65 \times 25$ making the answer $\boxed{\textbf{(C) 1,625}}$.

 \textbf{Remark} A good way to remember metric conversions is with the phrase \textbf{ "King Henry Died By Drinking Chocolate Milk" } (Kilo, Hecto, Deka, Basic units, Deci, Centi, Milli). As you read from left to right you divide by 10. So 1 meter (which is a basic unit) is equivalent to $1000$ millimeters, or $100$ centimeters. 

 


	\item (\textit{AMC 12 12B 2023}) You are playing a game. A $2 \times 1$ rectangle covers two adjacent squares (oriented either horizontally or vertically) of a $3 \times 3$ grid of squares, but you are not told which two squares are covered. Your goal is to find at least one square that is covered by the rectangle. A "turn" consists of you guessing a square, after which you are told whether that square is covered by the hidden rectangle. What is the minimum number of turns you need to ensure that at least one of your guessed squares is covered by the rectangle?


 $\textbf{(A)}~3\qquad\textbf{(B)}~5\qquad\textbf{(C)}~4\qquad\textbf{(D)}~8\qquad\textbf{(E)}~6$

 \textbf{Answer: }C



\textbf{Solution 1}

Notice that the $3\times3$ square grid has a total of $12$ possible $2\times1$ rectangles.

 Suppose you choose the middle square for one of your turns. The middle square is covered by $4$ rectangles, and each of the remaining $8$ squares is covered by a maximum of $2$ uncounted rectangles. This means that the number of turns is at least $1+\frac{12-4}{2}=1+4=5$.

 Now suppose you don't choose the middle square. The squares on the middle of the sides are covered by at most 3 uncounted rectangles, and the squares on the corners are covered by at most 2 uncounted rectangles. In this case, we see that the least number of turns needed to account for all 12 rectangles is $12\div 3=4.$ To prove that choosing only side squares indeed does cover all 12 rectangles, we need to show that the 3 rectangles per square that cover each side square do not overlap. Drawing the rectangles that cover one square, we see they form a $T$ shape and they do not cover any other side square. Hence, our answer is $4.$

                    \begin{figure}[H]
                        \centering
                        \includegraphics[width=0.5\textwidth]{Solution/AMC_12_12B_2023_5_Solution_Figure_1.png}
                    \end{figure}
                    




\textbf{Solution 2}

First, note that since the rectangle covers 2 squares, we only need to guess squares that are not adjacent to any of our other guesses. To minimize the amount of guesses, each of our guessed squares should try to touch another guess on one vertex and one vertex only. There are only two ways to do this: one with $5$ guesses, and one with $4$. Since the problem is asking for the minimum number, the answer is $\boxed{\text{(C) }   4}$.

 



\textbf{Solution 3}

Since the hidden rectangle can only hide two adjacent squares, we may think that we eliminate 8 squares and we're done, but think again. This is the AMC 10, so there must be a better solution (also note that every other solution choice is below 8 so we're probably not done) So, we think again, we notice that we haven't used the adjacent condition, and then it clicks. If we eliminate the four squares with only one edge on the boundary of the 9x9 square. We are left with 5 diagonal squares, since our rectangle can't be diagonal, we can ensure that we find it in 4 moves. So our answer is :$\boxed{\text{(C) }   4}$




\textbf{Solution 4 (Checkerboard)}

The $3 \times 3$ grid can be colored like a checkerboard with alternating black and white squares.Let the top left square be white, and the rest of the squares alternate colors, as shown below:
                    \begin{figure}[H]
                        \centering
                        \includegraphics[width=0.5\textwidth]{Solution/AMC_12_12B_2023_5_Solution_Figure_2.png}
                    \end{figure}
                    Each $2 \times 1$ rectangle always covers $1$ white square and $1$ black square.You can ensure that at least one of your guessed squares is covered by the rectangle by choosing either each of the white squares ($5$ turns) or each of the black squares ($4$ turns).

 Since it is ideal to be the most efficient with our turns, by choosing all the black squares, we guarantee that one of the $\boxed{\textbf{(C)}~4}$ squares are of the $2 \times 1$ rectangle. 

 (Minor edits by NSAoPS)

 



\textbf{Solution 5 (Logic)}

We realize that every $2 \times 1$ rectangle must contain an edge and no more than one edge. There are a total of four edges so the answer is $\boxed{\text{(C) }   4.}$.~darrenn.cp

 



\textbf{Solution 6 (Very Fast)}

Draw the 3x3 grid. We want to minimize our options, so we look at the square shared by the most 2x1 rectangles. This is the center square. So, to guarantee our success, we choose the four edges from the center square, so our answer is $\boxed{C.4}$



	\item (\textit{AMC 12 12B 2023}) When the roots of the polynomial


 
 \[P(x)  = (x-1)^1 (x-2)^2 (x-3)^3 \cdot \cdot \cdot (x-10)^{10}\]


 are removed from the number line, what remains is the union of 11 disjoint open intervals. On how many of these intervals is $P(x)$ positive?


 $\textbf{(A)}~3\qquad\textbf{(B)}~7\qquad\textbf{(C)}~6\qquad\textbf{(D)}~4\qquad\textbf{(E)}~5$

 \textbf{Answer: }C



\textbf{Solution 1}

The expressions to the power of even powers are always positive, so we don't need to care about those. We only need to care about $(x-1)^1(x-3)^3(x-5)^5(x-7)^7(x-9)^9$. We need 0, 2, or 4 of the expressions to be negative. The 9 through 10 interval and 10 plus interval make all of the expressions positive. The 5 through 6 and 6 through 7 intervals make two of the expressions negative. The 1 through 2 and 2 through 3 intervals make four of the expressions negative. There are $\boxed{\textbf{(C) 6}}$ intervals.

 



\textbf{Solution 2}

The roots of the factorized polynomial are intervals from numbers 1 to 10. We take each interval as being defined as the number behind it. To make the function positive, we need to have an even number of negative expressions. Real numbers raised to even powers are always positive, so we only focus on $x-1$, $x-3$, $x-5$, $x-7$, and $x-9$. The intervals 1 and 2 leave 4 negative expressions, so they are counted. The same goes for intervals 5, 6, 9, and 10. Intervals 3 and 4 leave 3 negative expressions and intervals 7 and 8 leave 1 negative expression. The solution is the number of intervals which is $\boxed{\textbf{(C) 6}}$.

 



\textbf{Solution 3}

We can use the turning point behavior at the roots of a polynomial graph to find out the amount of intervals that are positive.

 First, we evaluate any value on the interval $(-\infty, 1)$. Since the degree of $P(x)$ is $1+2+...+9+10$ = $\frac{10\times11}{2}$ = $55$, and every term in $P(x)$ is negative, multiplying $55$ negatives gives a negative value. So $(-\infty, 0)$ is a negative interval.

 We know that the roots of $P(x)$ are at $1,2,...,10$. When the degree of the term of each root is odd, the graph of $P(x)$ will pass through the graph and change signs, and vice versa. So at $x=1$, the graph will change signs; at $x=2$, the graph will not, and so on.

 This tells us that the interval $(1,2)$ is positive, $(2,3)$ is also positive, $(3,4)$ is negative, $(4,5)$ is also negative, and so on, with the pattern being $+,+,-,-,+,+,-,-,...$ .

 The positive intervals are therefore $(1,2)$, $(2,3)$, $(5,6)$, $(6,7)$, $(9,10)$, and $(10,\infty)$, for a total of $\boxed{\textbf{(C) 6}}$.

 



\textbf{Solution 4}

Denote by $I_k$ the interval $\left( k - 1 , k \right)$ for $k \in \left\{ 2, 3, \cdots , 10 \right\}$ and $I_1$ the interval $\left( - \infty, 1 \right)$.

 Therefore, the number of intervals that $P(x)$ is positive is
 \begin{align*} 1 + \sum_{i=1}^{10} \Bbb I \left\{  \sum_{j=i}^{10} j \mbox{ is even}   \right\}   & = 1 + \sum_{i=1}^{10} \Bbb I \left\{ \frac{\left( i + 10 \right) \left( 11 - i \right)}{2} \mbox{ is even}   \right\} \\  & = 1 + \sum_{i=1}^{10} \Bbb I \left\{ \frac{- i^2 + i + 110}{2} \mbox{ is even}   \right\} \\  & = 1 + \sum_{i=1}^{10} \Bbb I \left\{ \frac{i^2 - i}{2} \mbox{ is odd}   \right\} \\  & = \boxed{\textbf{(C) 6}} . \end{align*}





\textbf{Solution 5: Graphing}

Recall two key facts about the roots of a polynomial:- If a root has an odd multiplicity (e.g. the root appears an odd number of times), then the graph will cross the x-axis.- If a root has an even multiplicity (e.g. the root appears an even number of times), then the graph will bounce off the x-axis.

 Sketching the graph and noting the multiplicity of the roots, we see that there are $\boxed{C) 6}$ positive intervals.

 



\textbf{Solution 6: Solution 3 with powers}

When the range of \( P(x) \) is from \( -\infty \) to \( 1 \), the sign is negative because in  

 $(x - 1)^1 (x - 2)^2 (x - 3)^3 \cdots (x - 10)^{10}$,the terms corresponding to \( x = 1, 3, 5, 7, \) and \( 9 \) are all negative, and since they have odd exponents, they contribute a negative sign.  The other terms with even exponents (like \( (x - 2)^2 \), \( (x - 4)^4 \), etc.) are positive regardless of the sign inside, so they don't affect the overall sign. However, in the range\( \ 1 \) to \( \ 2 \), the power of $1$ flips, causing this range to be positive. For the range $2$ - $3$, the sign of any even inside the exponent doesn't matter, so this range stays positive. Continueing this pattern, We see the pattern is $+, +, -, -,$ and so on. there are 3 complete cycles, so $3 * 2 =$ $\boxed{\textbf{(C) 6}}$~BlueKite

 


	\item (\textit{AMC 12 12B 2023}) For how many integers $n$ does the expression
 \[\sqrt{\frac{\log (n^2) - (\log n)^2}{\log n - 3}}\]represent a real number, where log denotes the base $10$ logarithm?


 $\textbf{(A) }900 \qquad \textbf{(B) }3\qquad \textbf{(C) }902 \qquad \textbf{(D) } 2  \qquad \textbf{(E) }901$

 \textbf{Answer: }E



\textbf{Solution}

We have
 \begin{align*} \sqrt{\frac{\log \left( n^2 \right) - \left( \log n \right)^2}{\log n - 3}} & = \sqrt{\frac{2 \log n - \left( \log n \right)^2}{\log n - 3}} \\ & = \sqrt{\frac{\left( \log n \right) \left( 2 - \log n\right)}{\log n - 3}} . \end{align*}


 Because $n$ is an integer and $\log n$ is well defined, $n$ must be a positive integer.

 Case 1: $n = 1$ or $10^2$.

 The above expression is 0. So these are valid solutions.

 Case 2: $n \neq 1, 10^2$.

 Thus, $\log n > 0$ and $2 - \log n \neq 0$.To make the above expression real, we must have $2 < \log n < 3$.Thus, $100 < n < 1000$.Thus, $101 \leq n \leq 999$.Hence, the number of solutions in this case is 899.

 Putting all cases together, the total number of solutions is $\boxed{\textbf{(E) 901}}$.

 



\textbf{Solution (Solution 1 Simplified)}

Notice $\log(n^2)$ can be written as $2\log(n)$. Setting $a=\log(n)$, the equation becomes $\sqrt{\frac{2a-a^2}{a-3}}$ which can be written as $\sqrt{\frac{a(2-a)}{a-3}}$

 Case 1: $a \ge 3$The expression is undefined when $a=3$. For $a>3$, it is trivial to see that the denominator is positive and the numerator is negative, thus resulting in no real solutions.

 Case 2: $2 \le a<3$For $a=2$, the numerator is zero, giving us a valid solution. When $a>2$, both the denominator and numerator are negative so all real values of a in this interval is a solution to the equation. All integers of n that makes this true are between $10^2$ and $10^3-1$. There are 900 solutions here. 

 Case 3: $0<a<2$The numerator will be positive but the denominator is negative, no real solutions exist. 

 Case 4: $a=0$The expression evaluates to zero, $1$ valid solution exists. 

 Case 5: $a<0$All values for $a<0$ requires $0<n<1$, no integer solutions exist. 

 Adding up the cases:$900+1=\boxed{\textbf{(E) 901}}$




\textbf{Solution 3 (3 degree polynominal graph)}

for $\frac{a(2-a)}{a-3} >0$ , 

 transform it into a(a-2)(a-3) <0 , 

 use the following graph to quickly confirm 

 1) a < 0 or

 2) 2 < a <3

 
                    \begin{figure}[H]
                        \centering
                        \includegraphics[width=0.5\textwidth]{Solution/AMC_12_12B_2023_7_Solution_Figure_1.PNG}
                    \end{figure}
                    


	\item (\textit{AMC 12 12B 2023}) How many nonempty subsets $B$ of $\{0, 1, 2, 3, \dots, 12\}$ have the property that the number of elements in $B$ is equal to the least element of $B$? For example, $B = \{4, 6, 8, 11\}$ satisfies the condition.


 $\textbf{(A)}\ 256 \qquad\textbf{(B)}\ 136 \qquad\textbf{(C)}\ 108 \qquad\textbf{(D)}\ 144 \qquad\textbf{(E)}\ 156$

 \textbf{Answer: }D



\textbf{Solution}

There is no way to have a set with 0. If a set is to have its lowest element as 1, it must have only 1 element: 1. If a set is to have its lowest element as 2, it must have 2, and the other element will be chosen from the natural numbers between 3 and 12, inclusive. To calculate this, we do $\binom{10}{1}$. If the set is the have its lowest element as 3, the other 2 elements will be chosen from the natural numbers between 4 and 12, inclusive. To calculate this, we do $\binom{9}{2}$. We can see a pattern emerge, where the top is decreasing by 1 and the bottom is increasing by 1. In other words, we have to add $1 + \binom{10}{1} + \binom{9}{2} + \binom{8}{3} + \binom{7}{4} + \binom{6}{5}$. This is $1+10+36+56+35+6 = \boxed{\textbf{(D) 144}}$.

 \textbf{Note:} In general, i.e. when the number is not $13$, the answer is $F_{n-1}$.-Mr Sharkman

 


	\item (\textit{AMC 12 12B 2023}) What is the area of the region in the coordinate plane defined by


 $\left||x|-1\right|+\left||y|-1\right|\leq 1?$
$\textbf{(A)}~2\qquad\textbf{(B)}~8\qquad\textbf{(C)}~4\qquad\textbf{(D)}~15\qquad\textbf{(E)}~12$

 \textbf{Answer: }B



\textbf{Solution 1}

First consider, $|x-1|+|y-1| \le 1.$We can see that it is a square with a side length of $\sqrt{2}$ (diagonal $2$). The area of the square is $\sqrt{2}^2 = 2.$

 Next, we insert an absolute value sign into the equation and get $|x-1|+||y|-1| \le 1.$ This will double the square reflecting over x-axis.

 So now we have $2$ squares.

 Finally, we add one more absolute value and obtain $||x|-1|+||y|-1| \le 1.$ This will double the squares as we reflect the $2$ squares we already have over the y-axis.

 Concluding, we have $4$ congruent squares. Thus, the total area is $4\cdot2 =$ $\boxed{\text{(B) 8}}$




\textbf{Solution 2 (Graphing)}

We first consider the lattice points that satisfy $||x|-1| = 0$ and $||y|-1| = 1$. The lattice points satisfying these equationsare $(1,0), (1,2), (1,-2), (-1,0), (-1,2),$ and $(-1,-2).$ By symmetry, we also have points $(0,1), (2,1), (-2,1), (0,-1), (2,-1),$ and  $(-2,-1)$  when  $||x|-1| = 1$ and $||y|-1| = 0$.  Graphing and connecting these points, we form 5 squares. However,we can see that any point within the square in the middle does not satisfy the given inequality (take $(0,0)$, for instance). Asnoted in the above solution, each square has a diagonal $2$ for an area of $\frac{2^2}{2} = 2$, so the total area is $4\cdot2 =$$\boxed{\text{(B) 8}}.$ 

 \textbf{Note} This problem is very similar to a past AIME problem (1997 P13)

 \href{https://artofproblemsolving.com/wiki/index.php/1997\_AIME\_Problems/Problem\_13}{https://artofproblemsolving.com/wiki/index.php/1997\_AIME\_Problems/Problem\_13}

 



\textbf{Solution 3 (Logic)}

The value of $|x|$ and $|y|$ can be a maximum of 1 when the other is 0. Therefore the value of $x$ and $y$ range from -2 to 2. This forms a diamond shape which has area $4 \times \frac{2^2}{2}$ which is $\boxed{\text{(B) 8}}.$




\textbf{Solution 4}

We start by considering the graph of $|x|+|y|\leq 1$. To get from this graph to $||x|-1|+||y|-1| \leq 1$ we have to translate it by $\pm 1$ on the $x$ axis and $\pm 1$ on the $y$ axis.

 Graphing $|x|+|y|\leq 1$ we get a square with side length of $\sqrt{2}$, so the area of one of these squares is just $2$.

 We have to multiply by $4$ since there are $4$ combinations of shifting the $x$ and $y$ axis.

 So we have $2\times 4$ which is $\boxed{\text{(B) 8}}$.

 



\textbf{Solution 5 (Desperate)}

There are $2$ sets of $2$ absolute value bars. Each of those $2$ absolute value bars can take on $2$ values, so we have $2 \cdot 2 \cdot 2 = 8$ cases. We guess that the answer is divisible by $8$. The only answer choice that is divisible by $8$ is $\boxed{\text{(B)}~8}$.

 


	\item (\textit{AMC 12 12B 2023}) In the $xy$-plane, a circle of radius $4$ with center on the positive $x$-axis is tangent to the $y$-axis at the origin, and a circle with radius $10$ with center on the positive $y$-axis is tangent to the $x$-axis at the origin. What is the slope of the line passing through the two points at which these circles intersect?


 $\textbf{(A)}\ \dfrac{2}{7} \qquad\textbf{(B)}\ \dfrac{3}{7}  \qquad\textbf{(C)}\ \dfrac{2}{\sqrt{29}}  \qquad\textbf{(D)}\ \dfrac{1}{\sqrt{29}}  \qquad\textbf{(E)}\ \dfrac{2}{5}$

 \textbf{Answer: }E



\textbf{Solution 1}

The center of the first circle is $(4,0)$.The center of the second circle is $(0,10)$.Thus, the slope of the line that passes through these two centers is $- \frac{10}{4} = - \frac{5}{2}$.

 Because this line is the perpendicular bisector of the line that passes through two intersecting points of two circles, the slope of the latter line is $\frac{-1}{- \frac{5}{2}} = \boxed{\textbf{(E) } \frac{2}{5}}$.

 



\textbf{Solution 2 (Coordinate Geometry)}

The first circle can be written as $(x-4)^2 + y^2 = 4^2$ we'll call this equation $(1)$The second can we writen as $x^2 + (y-10)^2 = 10^2$, we'll call this equation $(2)$

 Expanding $(1)$:
 \begin{align*} x^2 -8x + 4^2 + y^2 &= 4^2 \\ x^2 - 8x + y^2 &= 0 \end{align*}
Exapnding $(2)$
 \begin{align*} x^2 + y^2 -20y + 10^2 = 10^2\\ x^2 + y^2 - 20y = 0 \end{align*}


 Now we can set the equations equal to eachother:
 \begin{align*} x^2 - 8x + y^2 &= x^2 + y^2 - 20y \\ \frac{8}{20}x &= y \\ \frac{2}{5}x &= y \end{align*}
This is in slope intercept form therefore the slope is $\boxed{\textbf{(E) } \frac{2}{5}}$.

 
                    \begin{figure}[H]
                        \centering
                        \includegraphics[width=0.5\textwidth]{Solution/AMC_12_12B_2023_10_Solution_Figure_1.PNG}
                    \end{figure}
                    



\textbf{Solution 3 (Polar Coordinates)}

In polar coordinates, a circle of radius $r$ tangent to the $y$-axis with center on the positive $x$-axis has equation $r = 2a\cos\theta$, and a circle tangent to the $x$-axis with center on the positive $y$-axis has equation $r = 2b\sin\theta$.

 The first circle (radius $4$, center $(4,0)$) has polar equation $r = 8\cos\theta$, and the second circle (radius $10$, center $(0,10)$) has polar equation $r = 20\sin\theta$.

 The origin lies on both circles, so it is one intersection point. Setting the equations equal for the other:
 \[8\cos\theta = 20\sin\theta \implies \tan\theta = \frac{8}{20} = \frac{2}{5}\]

 Since $\tan\theta = \frac{y}{x}$, both intersection points (the origin and the point at angle $\theta$) lie on the line through the origin with slope $\tan\theta$. Therefore, the slope of the line through the two intersection points is $\boxed{\dfrac{2}{5}}$.

 


	\item (\textit{AMC 12 12B 2023}) What is the maximum area of an isosceles trapezoid that has legs of length $1$ and one base twice as long as the other?


 $\textbf{(A) }\frac 54 \qquad \textbf{(B) } \frac 87 \qquad \textbf{(C)} \frac{5\sqrt2}4 \qquad \textbf{(D) } \frac 32  \qquad \textbf{(E) } \frac{3\sqrt3}4$

 \textbf{Answer: }D



\textbf{Solution 1}

Let the trapezoid be $ABCD$ with $AD = BC = 1, \; AB = x, CD = 2x$. Extend $AD$ and $BC$ to meet at point $E$. Then, notice $\triangle ABE \sim \triangle DCE$ with side length ratio $1:2$ and $AE = BE = 1$. Thus, $[DCE] = 4 \cdot [ABE]$ and $[ABCD] = [DCE] - [ABE] = \frac{3}{4} \cdot [DCE]$.

 The problem reduces to maximizing the area of $[DCE]$, an isosceles triangle with legs of length $2$. Analyzing the sine area formula, this is clearly maximized when $\angle DEC = 90^{\circ}$, so $[DCE] = 2$ and $[ABCD] = \frac{3}{4} \cdot 2 = \boxed{\frac{3}{2}}.$

 -PIDay

 



\textbf{Solution 2}

Denote by $x$ the length of the shorter base.Thus, the height of the trapezoid is
 \begin{align*} \sqrt{1^2 - \left( \frac{x}{2} \right)^2} . \end{align*}


 Thus, the area of the trapezoid is
 \begin{align*} \frac{1}{2} \left( x + 2 x \right) \sqrt{1^2 - \left( \frac{x}{2} \right)^2}  & = \frac{3}{4} \sqrt{x^2 \left( 4 - x^2 \right)} \\ & \leq \frac{3}{4} \frac{x^2 + \left( 4 - x^2 \right)}{2} \\ & = \boxed{\textbf{(D) } \frac{3}{2}} , \end{align*}


 where the inequality follows from the AM-GM inequality and it is binding if and only if $x^2 = 4 - x^2$.

 



\textbf{Solution 3 (Calculus)}

Derive the expression for area
 \[A = \frac{3}{4}x\sqrt{4-x^2}\]as in the solution above. To find the minimum, we can take the derivative with respect to $x$:
 \[\frac{dA}{dx} = \frac{3}{4}\sqrt{4-x^2}+\left(\frac{3x}{4}\right)\frac{-2x}{2\sqrt{4-x^2}} = \frac{6-3x^2}{2\sqrt{4-x^2}}.\]This expression is equal to zero when $x=\pm\sqrt{2}$, so $A$ has two critical points at $\pm\sqrt{2}$. But given the bounds of the problem, we can conclude $x = \sqrt{2}$ maximizes $A$ (alternatively you can do first derivative test). Plugging that value back in, we get $A_{\text{max}} = \boxed{(\text{D})\ \frac{3}{2}}$.

 (Slightly Simpler)

 Or rewrite the expression for area to be

 
 \[A = \frac{3}{4}\sqrt{4x^2-x^4}\]

 Now to find the minimum, we can just find the minimum of what's inside the square root (since the square root function is increasing). Take the derivative of $4x^2-x^4$,
 \[f'(x)=8x-4x^3.\]This is equal to zero at $x=0,\pm\sqrt{2}$ but the solution must be positive so $x=\sqrt{2}$.

 



\textbf{Solution 4 (Trigonometry)}

Let the length of the shorter base of the trapezoid be $2x$ and the height of the trapezoid be $y$.

 
                    \begin{figure}[H]
                        \centering
                        \includegraphics[width=0.5\textwidth]{Solution/AMC_12_12B_2023_11_Solution_Figure_1.png}
                    \end{figure}
                    

 Each leg has length $1$ if and only if $x^2+y^2=1$, where $x$ and $y$ are positive real numbers. The general solution to this equation is 
 \[(x,y)=(\cos t,\sin t)\] for any number $0<t<\frac{\pi}{2}$ so that $x$ and $y$ are positive. The area to maximize is 
 \[\frac{1}{2}(2x+4x)y=3xy\] Hence, we maximize $3\sin t\cos t$ for $0<t<\frac{\pi}{2}$.
 \begin{align*}     3xy &= 3\sin t\cos t \\     &= \frac{3}{2}(2\sin t\cos t) \\     &= \frac{3}{2}\sin(2t) \end{align*}
The maximum of $\sin(2t)$ is $1$, thus the maximum of $3xy$ is $\boxed{\text{(D) }\frac{3}{2}}$ which occurs at $t=\frac{\pi}{4}$, satisfying the inequality $0<t<\frac{\pi}{2}$.

 


 



\textbf{Solution 5}

Denote $x$ and $2x$ as the two bases. We can straightforwardly find that the height of the trapezoid is $\sqrt{1-\frac{x^2}{4}}$ by the Pythagorean theorem. 

 The area of this trapezoid is then given by the expression $\frac{x+2x}{2}\cdot\sqrt{1-\frac{x^2}{4}}.$

 Let $m$ be the maximum value that this expression achieves. Since $x$ is positive and assuming $m\geq 1$, we can perform the following operations to maximize $16m^2$, thus $m^2$, and thus $m$. We have

 
 \[\frac{x+2x}{2}\cdot\sqrt{1-\frac{x^2}{4}} = m\]
 \[\frac{9x^2}{4}\cdot\frac{4-x^2}{4} = m^2\]
 \[-9x^4 + 36x^2 = 16m^2.\]

 The maximum value of this quartic occurs at $x^2 = \frac{-36}{2(-9)} = 2$, or when $x = \sqrt{2}.$

 It follows that the area of the trapezoid is equal to 

 
 \[\frac{3\sqrt{2}}{2}\cdot\frac{1}{\sqrt{2}} = \boxed{\text{(D) }\frac{3}{2}}.\]

 -Benedict T (countmath1)

 


	\item (\textit{AMC 12 12B 2023}) For complex numbers $u=a+bi$ and $v=c+di$, define the binary operation $\otimes$ by
 \[u\otimes v=ac+bdi.\]Suppose $z$ is a complex number such that $z\otimes z=z^{2}+40$. What is $|z|$?


 $\textbf{(A) }2\qquad\textbf{(B) }5\qquad\textbf{(C) }\sqrt{5}\qquad\textbf{(D) }\sqrt{10}\qquad\textbf{(E) }5\sqrt{2}$

 \textbf{Answer: }E



\textbf{Solution 1}

let $z$ = $a+bi$.

 $z \otimes z = a^{2}+b^{2}i$.

 This is equal to $z^{2} + 40 = a^{2}-b^{2}+40+2abi$

 Since the real values have to be equal to each other, $a^{2}-b^{2}+40 = a^{2}$.Simple algebra shows $b^{2} = 40$, so $b = \pm 2\sqrt{10}$.

 The imaginary components must also equal each other, meaning $b^{2} = 2ab$, or $b = 2a$. This means $a = \frac{b}{2} = \pm \sqrt{10}$.

 Thus, the magnitude of $z$ is $\sqrt{a^{2}+b^{2}} = \sqrt{10+40} = \sqrt{50} = 5\sqrt{2}$$=\text{\boxed{\textbf{(E) }5\sqrt{2}}}$



	\item (\textit{AMC 12 12B 2023}) A rectangular box $P$ has distinct edge lengths $a$, $b$, and $c$. The sum of the lengths of all $12$ edges of $P$ is $13$, the sum of the areas of all $6$ faces of $P$ is $\frac{11}{2}$, and the volume of $P$ is $\frac{1}{2}$. What is the length of the longest interior diagonal connecting two vertices of $P$?


 $\textbf{(A)}~2\qquad\textbf{(B)}~\frac{3}{8}\qquad\textbf{(C)}~\frac{9}{8}\qquad\textbf{(D)}~\frac{9}{4}\qquad\textbf{(E)}~\frac{3}{2}$

 \textbf{Answer: }D



\textbf{Solution 1 (algebraic manipulation)}


                    \begin{figure}[H]
                        \centering
                        \includegraphics[width=0.5\textwidth]{Solution/AMC_12_12B_2023_13_Solution_Figure_1.png}
                    \end{figure}
                    

 We can create three equations using the given information. 
 \[4a+4b+4c = 13\]
 \[2ab+2ac+2bc=\frac{11}{2}\]
 \[abc=\frac{1}{2}\]We also know that we want $\sqrt{a^2 + b^2 + c^2}$ because that is the length that can be found from using the Pythagorean Theorem. We cleverly notice that $a^2 + b^2 + c^2 = (a+b+c)^2 - 2(ab+ac+bc)$. We know that $a+b+c = \frac{13}{4}$ and $2(ab+ac+bc)=\dfrac{11}2$, so $a^2 + b^2 + c^2 = \left(\frac{13}{4}\right)^2 - \frac{11}{2} = \frac{169-88}{16} = \frac{81}{16}$. So our answer is $\sqrt{\frac{81}{16}} = \boxed{\textbf{(D)}~\frac94}$.

 



\textbf{Solution 2 (Vieta's)}

We use the equations from Solution 1 and manipulate it a little:
 \[a+b+c = \frac{13}{4}\]
 \[ab+ac+bc=\frac{11}{4}\]
 \[abc=\frac{1}{2}\]Notice how these are the equations for the vieta's formulas for a polynomial with roots of $a$, $b$, and $c$. Let's create that polynomial. It would be $x^3 - \frac{13}{4}x^2 + \frac{11}{4}x - \frac{1}{2}$. Multiplying each term by 4 to get rid of fractions, we get $4x^3 - 13x^2 + 11x - 2$. Notice how the coefficients add up to $0$. Whenever this happens, that means that $(x-1)$ is a factor and that 1 is a root. After using synthetic division to divide $4x^3 - 13x^2 + 11x - 2$ by $x-1$, we get $4x^2 - 9x + 2$. Factoring that, you get $(x-2)(4x-1)$. This means that this polynomial factors to $(x-1)(x-2)(4x-1)$ and that the roots are $1$, $2$, and $1/4$. Since we're looking for $\sqrt{a^2 + b^2 + c^2}$, this is equal to $\sqrt{1^2 + 2^2 + \frac{1}{4}^2} = \sqrt{\frac{81}{16}} = \boxed{\textbf{(D)}~\frac94}$




\textbf{Solution 3 (Newton's Sums)}

Same formal process of Solution 2, once we have $4x^3 - 13x^2 + 11x - 2$. Let $s_{2} = a^2 + b^2 + c^2$ and $s_{1} = a + b + c = \frac{13}{4}$. Applying \href{https://artofproblemsolving.com/wiki/index.php/Newton's\_Sums}{Newton's Sums}, we have 
 \[(4)s_{2} + (-13)s_{1} + 2 * 11 = 0\]. We can solve that $s_{2} = \frac{81}{16}$, and the interior diagonal would be 
 \[\sqrt{s_{2}} = \frac{9}{4}\].

 



\textbf{Solution 4 (Cheese Method)}

Incorporating the solution above, we know $a+b+c$ = $13/4$ $\Rightarrow$ $a+b+c > 3$. The side lengths are larger than $1$ $\cdot$ $1$ $\cdot$ $1$ (a unit cube). The side length of the interior of a unit cube is $\sqrt{3}$, and we know that the side lengths are larger than $1$ $\cdot$ $1$ $\cdot$ $1$, so that means the diagonal has to be larger than $\sqrt{3}$, and the only answer choice larger than $\sqrt{3}$ $\Rightarrow$ $\boxed{\textbf{(D)}~\frac94}$

 Note that the real number $\sqrt{3}$ is around $1.73$. Option $A$ is also greater than $\sqrt{3}$ meaning there are two options greater than $\sqrt{3}$. Option $A$ is an integer so educationally guessing we arrive at answer $D$ $\Rightarrow$ $\boxed{\textbf{(D)}~\frac94}$



	\item (\textit{AMC 12 12B 2023}) For how many ordered pairs $(a,b)$ of integers does the polynomial $x^3+ax^2+bx+6$ have $3$ distinct integer roots?


 $\textbf{(A)}\ 5 \qquad\textbf{(B)}\ 6 \qquad\textbf{(C)}\ 8 \qquad\textbf{(D)}\ 7 \qquad\textbf{(E)}\ 4$

 \textbf{Answer: }A



\textbf{Solution}

Denote three roots as $r_1 < r_2 < r_3$.Following from Vieta's formula, $r_1r_2r_3 = -6$.

 Case 1: All roots are negative.

 We have the following solution: $\left( -3, -2, -1 \right)$.

 Case 2: One root is negative and two roots are positive.

 We have the following solutions: $\left( -3, 1, 2 \right)$, $\left( -2, 1, 3 \right)$, $\left( -1, 2, 3 \right)$, $\left( -1, 1, 6 \right)$.

 Putting all cases together, the total number of solutions is$\boxed{\textbf{(A) 5}}$.

 


	\item (\textit{AMC 12 12B 2023}) Suppose $a$, $b$, and $c$ are positive integers such that
 \[\frac{a}{14}+\frac{b}{15}=\frac{c}{210}.\]Which of the following statements are necessarily true?


 I. If $\gcd(a,14)=1$ or $\gcd(b,15)=1$ or both, then $\gcd(c,210)=1$.


 II. If $\gcd(c,210)=1$, then $\gcd(a,14)=1$ or $\gcd(b,15)=1$ or both.


 III. $\gcd(c,210)=1$ if and only if $\gcd(a,14)=\gcd(b,15)=1$.


 $\textbf{(A)}~\text{I, II, and III}\qquad\textbf{(B)}~\text{I only}\qquad\textbf{(C)}~\text{I and II only}\qquad\textbf{(D)}~\text{III only}\qquad\textbf{(E)}~\text{II and III only}$

 \textbf{Answer: }E



\textbf{Solution 1 (Guess and check + Contrapositive)}

We examine each of the conditions. 

 The first condition is false. A simple counterexample is $a=3$ and $b=5$. The corresponding value of $c$ is $115$. Since $\gcd(3,14)=1$, condition $I$ would imply that $\gcd(c,210)=1.$ However, $\gcd(115,210)$ is clearly not $1$ (they share a common factor of $5$). Condition $I$ is false so that we can rule out choices $A,B,$ and $C$. 

 We now decide between the two answer choices $D$ and $E$. What differs between them is the validity of condition $II$, so it suffices to check $II$ simply. 

 We look at statement $II$'s contrapositive to prove it. The contrapositive states that if $\gcd(a,14)\neq1$ and $\gcd(b,15)\neq1$, then $\gcd(c,210)\neq1.$ In other words, if $a$ shares some common factor that is not $1$ with $14$ and $b$ shares some common factor that is not $1$ with $15$, then $c$ also shares a common factor that is not $1$ with $210$. Let's say that $a=a'\cdot n$, where $a'$ is a factor of $14$ not equal to $1$. (So $a'$ is the common factor.) 

 We can rewrite the given equation as $15a+14b=c\implies15(a'n)+14b=c.$ We can express $14$ as $a'\cdot n'$, for some positive integer $n'$ (this $n'$ can be $1$). We can factor $a'$ out to get $a'(15n+bn')=c.$ 

 Since all values in this equation are integers, $c$ must be divisible by $a'$. Since $a'$ is a factor of $14$, $a'$ must also be a factor of $210$, a multiple of $14$. Therefore, we know that $c$ shares a common factor with $210$ (which is $a'$), so $\gcd(c,210)\neq1$. This is what $II$ states, so therefore $II$ is true. 

 Thus, our answer is $\boxed{\textbf{(E) }\text{II and III only}}.$ ~Technodoggo~AVM2023 (Minor Edits - Fixes in the first paragraph and grammar edits)

 



\textbf{Solution 2}

The equation given in the problem can be written as
 \[ 15 a + 14 b = c. \hspace{1cm} (1) \]
$\textbf{First, we prove that Statement I is not correct.}$

 A counter example is $a = 1$ and $b = 3$.Thus, ${\textrm gcd} (c, 210) = 3 \neq 1$.

 $\textbf{Second, we prove that Statement III is correct.}$

 First, we prove the ``if\textit{ part. }

 Suppose ${\textrm gcd}(a , 14) = 1$ and ${\textrm gcd}(b, 15) = 1$. However, ${\textrm gcd} (c, 210) \neq 1$.

 Thus, $c$ must be divisible by at least one factor of 210. W.L.O.G, we assume $c$ is divisible by 2.

 Modulo 2 on Equation (1), we get that $2 | a$.This is a contradiction with the condition that ${\textrm gcd}(a , 14) = 1$.Therefore, the ``if\textit{ part in Statement III is correct.}

 Second, we prove the ``only if\textit{ part.}

 Suppose ${\textrm gcd} (c, 210) \neq 1$. Because $210 = 14 \cdot 15$, there must be one factor of 14 or 15 that divides $c$.W.L.O.G, we assume there is a factor $q > 1$ of 14 that divides $c$.Because ${\textrm gcd} (14, 15) = 1$, we have ${\textrm gcd} (q, 15) = 1$.Modulo $q$ on Equation (1), we have $q | a$.

 Because $q | 14$, we have ${\textrm gcd}(a , 14) \geq q > 1$.

 Analogously, we can prove that ${\textrm gcd}(b , 15) > 1$.

 $\textbf{Third, we prove that Statement II is correct.}$

 This is simply a special case of the ``only if\textit{ part of Statement III. So we omit the proof.}

 All analyses above imply$\boxed{\textbf{(E) II and III only}}.$




\textbf{Solution 3 (Answer Choices)}

It can easily be shown that statement I is false (a counterexample would be $a=1, b=5, c=85$), meaning the only viable answer choices are D and E. Since both of these answer choices include statement III, this means III is true. Since III is true, this actually implies that statement II is true, as III is just a stronger version of statement II (or its contrapositive, to be precise). Therefore the answer is $\boxed{\textbf{(E) II and III only}}.$




\textbf{Solution 4}

First, we will show that statement I is false. A simple counterexample is $a=2$ and $b=2$. Then, $c$ is even, and it is not relatively prime with $210$.

 
Now, we will focus on statement III. It's pretty easy to realize that if either $a$ isn't relatively prime to $14$ or $b$ isn't relatively prime to $15$, it will always share factors with $210$ because $c = 15a + 14b$ (think $15 \cdot 14 + 14$). Therefore, statement III must be true. 

 
Since statement III is a stricter case of statement II, the last two statements are true, hence choose $\boxed{\textbf{(E) }\text{II and III only}}.$




\textbf{Solution 5 (Rigorous)}

The condition in the problem is equivalent to\begin{align*}15 a + 14 b = c. \hspace{1cm} \end{align*}


 From here we can use the properties of the greatest common divisor to get

 
 \[\gcd(c, 14)=\gcd(15a+14b, 14)=\gcd(15a, 14)=\gcd(a, 14)\]

 and

 
 \[\gcd(c, 15)=\gcd(15a+14b, 15)=\gcd(14b, 15)=\gcd(b, 15)\]

 From here it's not hard to get

 
 \[\gcd(c, 210)=\gcd(c, 15)*\gcd(c, 14)=\gcd(a, 14)*\gcd(b, 15)\]

 Hence $\gcd(c, 210)=1$ is equivalent to $\gcd(a, 14)=1$ and $\gcd(b, 15)=1$. This means the answer must be $\boxed{\textbf{(E) }\text{II and III only}}.$




\textbf{Solution 6 (Intuitive Risking)}

We know that \( \frac{a}{14} + \frac{b}{15} = \frac{c}{210} \). So, therefore, \( \frac{a + 14b}{15} = \frac{c}{1} \), and \( 15a + 14b = 210c \).

 We could see that some value of \(a\) and another value of \(b\) gives a different value of \(c\). \href{https://artofproblemsolving.com/wiki/index.php?title=WLOG}{WLOG}, we give an intuitive risk, saying that \( a + b = c \). Our goal is to just prove statements, not find solutions.

 Statement I:

 If \( \gcd(a,14) = 1 \) or \( \gcd(b,15) = 1 \) or both, then \( \gcd(c,210) = 1 \).

 We have \( a + b = c \).

 Then if \( \gcd(a,14) = 1 \), we say \( a = 1 \)

 So then \( \gcd(b,14) \neq 1 \), so we say \( b = 2 \).

 \( 1 + 2 = 3 \), and \( \gcd(3,210) = 1 \) is false. So statement I is incorrect.

 Statement II:

 If \( \gcd(c,210) = 1 \), then \( \gcd(a,14) = 1 \) or \( \gcd(b,15) = 1 \) or both.

 We say  \( a + b = c \).

 For statement II, we need to prove that it works for every odd number that is a multiple of the first number that makes both \(a\) and \(b\) positive integers and \( \gcd(c,210) = 1 \). 

 This number is 9. The possible pair values for \(a\) and \(b\) are 5 and 4 respectively. If we try 27, we get for a and b can be 15 and 12. We clearly see this is a pattern, and to get to an odd multiple of 9 we just multiply the sums by the \( n \)th multiple of 9. Therefore statement II is always true.

 Statement III:

 \( \gcd(c,210) = 1 \) if and only if \( \gcd(a,14) = \gcd(b,15) = 1 \).

 We say \( a + b = c \)

 We prove this by contradiction. 

 Assume \( \gcd(a,14) \neq \gcd(b,15) \neq 1 \). 

 Then, that means \( \gcd(c,210) \) is never equal to 1. 

 So if we take \( a = 14 \) and \( b = 15 \), we get that \( \gcd(29,210) \neq 1 \). This is obviously false, because 29 is prime. 

 Therefore statement III is also correct.

 Only II and III are correct, leading us to answer choice $\boxed{\textbf{(E) }\text{II and III only}}.$




\textbf{Solution 7 (Formal Logic Lattice)}

Multiply the original equation by $210$ to get $15a+14b=c$.

 Denote the propositions $A:=(\gcd(a,14)=1)$, $B:=(\gcd(b,15)=1)$, $C:=(\gcd(c,210)=1)$.

 Work on the lattice with set $\{ A,B,C \}$ with $\ge$ the order relation given by $\implies$, $\land$ the logical minimum (AND), $\lor$ the logical maximum (OR), $\equiv$ logical equivalence $\iff$.

 I. $A\lor B\ge C\iff (A\ge C)\land (B\ge C)$

 II. $C\ge A\lor B$

 III. $C\equiv A\land B$

 Note $A\land B\ge A\lor B$, so $\mathbf{III}\implies \mathbf{II}$. Also, $\mathbf{III}\land\mathbf{I}\implies A\equiv B\equiv C$ by the squeeze lemma. In our case $A\not\equiv B\implies\neg \mathbf{III}\lor \neg \mathbf{I}$ (at least one of these statements is false). (The squeeze lemma refers to $A\land B\ge A\lor B\implies A\equiv B$ from decomposing $\mathbf{I}$ into $A\ge C\equiv A\land B \iff A\ge B$ and symmetrically $B\ge A$.)

 But notice $\gcd(14,15)=1$. Use multiplicativity $\gcd(x,y)=1\implies \gcd(z,xy)=\gcd(z,x)\gcd(z,y)$ and Euclidean algorithm $\gcd(x,y)=\gcd(x,y\pm x)$ to get

 
 \[\gcd(c,210)=\gcd(15a+14b,14)\gcd(15a+14b,15)=\gcd(a,14)\gcd(b,15).\]

 Since $\gcd: \mathbb{Z}\times \mathbb{Z}\to \mathbb{Z}_{\ge 0} \implies \gcd\ge 1$, so $\gcd(c,210)=1\iff \gcd(a,14)=\gcd(b,15)=1$. By the above, $\mathbf{III}$ is true, $\mathbf{II}$ is true, $\mathbf{I}$ is false, so the answer is $\boxed{\textbf{(E) }\text{II and III only}}$.

 



\textbf{Solution 8 (Number Theory)}

Rewrite as $15a + 14b = c$. Assume I, particularly the "or both" section. this fails if $a = 1$ and $b = 2$, for which we have the $\gcd$ is a factor of 2 times et cetera. Due to failure within the "or both" case, note that we can let $a = 2$ for which the first half fails. I is false.

 Consider II. Firstly, we show "or both" is a viable case. Consider $\gcd(c,210) = 1$. Then, $\gcd(14a+15b,210) = 1$. So, $14a+15b$ is odd not a multiple of 3 nor 7. Modulo 3 gives $2a \equiv 0 \pmod{3}$. That implies $a \equiv 0 \pmod{3}$. Also, take modulo 7. Then, $b \equiv 0 \pmod{7}$. $\gcd(14a+15b, 210) > 1$ holds if $a,b$ is a multiple of either 7 or 3 respectively (or both). This implies that if $\gcd(14a+15b, 210) = 1$ holds, then none of $a,b$ is a multiple of 7 or 3, meaning that $\gcd(a,14) = \gcd(b,15) = 1$, which satisfies the "or both" case of II, and thus II is true.

 Consider III. That is literally the "both" case of II, which we proved is always true.

 The answer is $\boxed{\text{II and III only}}$.

 



\textbf{Solution 9 (Set Theory)}

We can use set theory to approach this problem, firstly by defining the universal set $U$ to be all possible positive integer ordered pairs $(a, b)$.

 Next, create a Venn diagram that consists of 2 sets: a set $A$ of pairs $(a, b)$ so that $\gcd(a, 14) > 1$, and a set $B$ of all pairs $(a, b)$ so that $\gcd(b, 15) > 1$. In other words, set $A$ includes all pairs where $a$ is a multiple of $2$ or $7$ and set $B$ includes all pairs where $b$ is a multiple of $3$ or $5$. 

 Multiplying the original equation by $210$ yields $15a + 14b = c$. Because $210 = 2 \cdot 3 \cdot 5 \cdot 7$, if $\gcd(c, 210) > 1$, $a$ must be a multiple of $2$ or $7$ or $b$ must be a multiple of $3$ or $5$. Thus, $\gcd(c, 210) > 1$ directly corresponds to $A \cup B$.

 Next, understand what it means when the $\gcd$ function equals $1$. $\gcd(a, 14) = 1$, for example, means that $a$ is $\textbf{not}$ a multiple of $2$ or $7$. Similarly, $\gcd(b, 15) = 1$ means that $b$ is $\textbf{not}$ a multiple of $3$ or $5$. Note that this means the set of pairs satisfying $\gcd(a, 14) = 1$ is the complement of set $A$ ($A'$). Similarly, the set of pairs satisfying $\gcd(b, 15) = 1$ is the complement of set $B$ ($B'$) and the set of pairs satisfying $\gcd(c, 210) = 1$ is the complement of $A \cup B$ ($(A \cup B)'$).

 Therefore, we can rewrite statements I, II, and III using the complements of set A, set B, and $A \cup B$.

 I. $A' \cup B' \subseteq (A \cup B)'$

 II. $(A \cup B)' \subseteq A' \cup B’$

 III. $(A \cup B)' = A' \cap B’$
$A’ \cup B’$, $(A \cup B)'$, and $A’ \cap B’$ can be modeled as follows:

 <table align="center" cellpadding="5"><tbody><tr><td>
                    \begin{figure}[H]
                        \centering
                        \includegraphics[width=0.5\textwidth]{Solution/AMC_12_12B_2023_15_Solution_Figure_1.png}
                    \end{figure}
                    
</td><td>
                    \begin{figure}[H]
                        \centering
                        \includegraphics[width=0.5\textwidth]{Solution/AMC_12_12B_2023_15_Solution_Figure_2.png}
                    \end{figure}
                    
</td><td>
                    \begin{figure}[H]
                        \centering
                        \includegraphics[width=0.5\textwidth]{Solution/AMC_12_12B_2023_15_Solution_Figure_3.png}
                    \end{figure}
                    
</td></tr></tbody></table>If we simply compare these Venn diagrams to the rewritten statements, we can see that I is false and II and III (which turns out to be De Morgan's First Law) are true. Thus, the answer is $\boxed{\text{(E) II and III only}}$.

 


	\item (\textit{AMC 12 12B 2023}) In Coinland, there are three types of coins, each worth $6, 10,$ and $15.$ What is the sum of the digits of the maximum amount of money that is impossible to have?


 $\textbf{(A) }8\qquad\textbf{(B) }10\qquad\textbf{(C) }7\qquad\textbf{(D) }11\qquad\textbf{(E) }9$

 \textbf{Answer: }D



\textbf{Solution 1}

This problem asks to find largest $x$ that cannot be written as
 \[6 a + 10 b + 15 c = x, \hspace{1cm} (1)\]

 where $a, b, c \in \Bbb Z_+$.

 Denote by $r \in \left\{ 0, 1 \right\}$ the remainder of $x$ divided by 2.Modulo 2 on Equation (1), we getBy using modulus $m \in \left\{ 2, 3, 5 \right\}$ on the equation above, we get$c \equiv r \pmod{2}$. 

 Following from the \href{https://artofproblemsolving.com/wiki/index.php?title=Chicken\_McNugget\_Theorem}{Chicken McNugget Theorem}, we have that any number that is no less than $(3-1)(5-1) = 8$ can be expressed in the form of $3a + 5b$ with $a, b \in \Bbb Z_+$.

 Therefore, all even numbers that are at least equal to $2 \cdot 8 + 15 \cdot 0 = 16$ can be written in the form of Equation (1) with $a, b, c \in \Bbb Z_+$.All odd numbers that are at least equal to $2 \cdot 8 + 15 \cdot 1 = 31$ can be written in the form of Equation (1) with $a, b, c \in \Bbb Z_+$.

 The above two cases jointly imply that all numbers that are at least 30 can be written in the form of Equation (1) with $a, b, c \in \Bbb Z_+$.

 Next, we need to prove that 29 cannot be written in the form of Equation (1) with $a, b, c \in \Bbb Z_+$.

 Because 29 is odd, we must have $c \equiv 1 \pmod{2}$.Because $a, b, c \in \Bbb Z_+$, we must have $c = 1$.Plugging this into Equation (1), we get $3 a + 5 b = 7$.However, this equation does not have non-negative integer solutions.

 All analysis above jointly imply that the largest $x$ that has no non-negative integer solution to Equation (1) is 29.Therefore, the answer is $2 + 9 = \boxed{\textbf{(D) 11}}$.

 



\textbf{Solution 2 (Multiples Of 6)}

Arrange the positive integers into rows of $6$, like this:
 \[\begin{array}{|c|c|c|c|c|c|} \hline 01 & 02 & 03 & 04 & 05 & \textcolor{red}{\boxed{06}} \\ \hline 07 & 08 & 09 & \textcolor{red}{\boxed{10}} & 11 & \textcolor{red}{\boxed{12}} \\ \hline 13 & 14 & \textcolor{red}{\boxed{15}} & \textcolor{red}{\boxed{16}} & 17 & \textcolor{red}{\boxed{18}} \\ \hline 19 & \textcolor{red}{\boxed{20}} & \textcolor{red}{\boxed{21}} & \textcolor{red}{\boxed{22}} & 23 & \textcolor{red}{\boxed{24}} \\ \hline \textcolor{red}{\boxed{25}} & \textcolor{red}{\boxed{26}} & \textcolor{red}{\boxed{27}} & \textcolor{red}{\boxed{28}} & 29 & \textcolor{red}{\boxed{30}} \\ \hline \textcolor{red}{\boxed{31}} & \textcolor{red}{\boxed{32}} & \textcolor{red}{\boxed{33}} & \textcolor{red}{\boxed{34}} & \textcolor{red}{\boxed{35}} & \textcolor{red}{\boxed{36}} \\ \hline \end{array}\]
 \[\vdots\]Observe that if any number can be made from a combination of $6$s, $10$s, and $15$s, then every number below it in the same column must also be possible to make, by simply adding 6s.

 Thus, we will cross out any numbers that CAN be made as well as all numbers below it.

 In column 1, $25$ is possible $(10+15)$ and so is everything below $25$.

 Column 2 - cross out $20$

 Column 3 - cross out $15$

 Column 4 - cross out $10$

 Column 5 - cross out $35$

 Column 6 - cross all out.

 The maximum number that remains is $29$.Answer is $2+9= \boxed{\textbf{(D) 11}}$.

 



\textbf{Solution 3 (Modulo 6)}

Let the number of $6$ cent coins be $a$, the number of $10$ cent coins be $b$, and the number of $15$ cent coins be $c$. We get the \href{https://artofproblemsolving.com/wiki/index.php?title=Diophantine\_equation}{Diophantine equation}

 
 \[6a + 10b + 15c = k\]

 and we wish to find the largest possible value of $k$

 Construct the following $\mod 6$ table of $6$, $10$, and $15$

 \[\begin{array}{c|ccc}  &  &  & \\ \text{number of coins} & 6 & 10 & 15 \\  \hline  &  &  & \\  1 & 0 & 4 & 3\\  &  &  &  \\ 2 & 0 & 2 & 0 \\ \end{array}\]

 There are only $6$ possible residues for $6$, they are: $0$, $1$, $2$, $3$, $4$, and $5$.

 \begin{verbatim}
$0 \mod 6$We can obtain any value that is $0 \mod 6$ because we have $6$ cent coins.
\end{verbatim}

\begin{verbatim}
$1 \mod 6$We can obtain values that equal $1 \mod 6$ by using one $10$ cent coin, one $15$ cent coin, and as many $6$ cent coins needed.The smallest value that equals $1 \mod 6$ we can obtain is $10+15 = 25$. Therefore, the largest value that equals $1 \mod 6$ we cannot obtain is $25-6=19$
\end{verbatim}

\begin{verbatim}
$2 \mod 6$We can obtain values that equal $2 \mod 6$ by using two $10$ cent coins, and as many $6$ cent coins as needed.The smallest value that equals $2 \mod 6$ we can obtain is $10+10 = 20$. Therefore, the largest value that equals $2 \mod 6$ we cannot obtain is $20-6=14$
\end{verbatim}

\begin{verbatim}
$3 \mod 6$We can obtain values that equal $3 \mod 6$ by using one $15$ cent coin, and as many $6$ cent coins as needed.The smallest value that equals $3 \mod 6$ we can obtain is $15$. Therefore, the largest value that equals $3 \mod 6$ we cannot obtain is $15-6=9$
\end{verbatim}

\begin{verbatim}
$4 \mod 6$We can obtain values that equal $4 \mod 6$ by using one $10$ cent coin, and as many $6$ cent coins as needed.The smallest value that equals $4 \mod 6$ we can obtain is $10$. Therefore, the largest value that equals $4 \mod 6$ we cannot obtain is $10-6=4$
\end{verbatim}

\begin{verbatim}
$5 \mod 6$We can obtain values that equal $5 \mod 6$ by using two $10$ cent coins, one $15$ cent coin, and as many $6$ cent coins as needed.The smallest value that equals $5 \mod 6$ we can obtain is $10+10+15 = 35$. Therefore, the largest value that equals $5 \mod 6$ we cannot obtain is $35-6=29$
\end{verbatim}

Hence, the largest value in cents we cannot obtain using $6$, $10$, and $15$ cent coins is $29$, $2 + 9 = \boxed{\textbf{(D) 11}}$.

 



\textbf{Solution 4}

We claim that the largest number that cannot be obtained using $6$, $10$, and $15$ cent coins is $29$.

 Let's first focus on the combination of $6$, $10$. As both of them are even numbers, we cannot obtain any odd numbers from these two but requires $15$ to sum up to an odd number. Notice that by Chicken McNugget Theorem, the largest even number cannot be obtained by $6$, $10$ is $2(3\cdot 5-3-5)=14$. Add this with $29$, we can easily verify that $29$ cannot be obtained by $6$, $10$, and $15$ as it needs at least one odd number, with the remaining part cannot be represented by $6$ and $10$.

 Let's show that any number greater than $29$ can be obtained. First, any even numbers greater than $29$ can be obtained by $6$ and $10$ by the Chicken McNugget Theorem. Next, any odd number greater than $29$ can be obtained by adding one $15$ with some $6$s and $10$s, which is also shown by the Chicken McNugget Theorem. This completes the proof. So the answer is $2+9 = \boxed{\textbf{(D) 11}}$




\textbf{Solution 5 (apply Chicken McNugget Theorem twice)}

First considering the two terms 

 
 \[6a+10b = 2(3a+5b)\]

 \[3a+5b = (3-1)(5-1)+t = 8+t, t \geq 0\]

 Again applying the two variable formula for the terms in brackets we see that

 
 \[6a+10b+15c = 2(8+t)+15c = 2t+15c+16 = (2-1)(15-1)+s+16 = 30+s, s>=0\]

 so the given expression 6a+10b+15c produces all numbers from 30 and 29 is the largest number that could not be produced, so the answer is $2+9 = \boxed{\textbf{(D) 11}}$

            \begin{itemize}
        \item Note: Chicken McNugget Theorem: any positive integer N= (a-1)(b-1)+t ( a,b co-prime, t>=0) can always be represented as N = aP+ bQ ( p,q are non-negative integer)
            \end{itemize}
            NOTE (by Mqoi\_KOLA)- this method isn't true in general for any three coprime numbers take for example $10,14,3$ from the method used here one get $49$ as the largest possible number which can't be achieved which is false because $11$ is unachievable.

 


	\item (\textit{AMC 12 12B 2023}) Triangle $ABC$ has side lengths in arithmetic progression, and the smallest side has length $6$. If the triangle has an angle of $120^\circ,$ what is the area of $ABC$?


 $\textbf{(A) }12\sqrt{3}\qquad\textbf{(B) }8\sqrt{6}\qquad\textbf{(C) }14\sqrt{2}\qquad\textbf{(D) }20\sqrt{2}\qquad\textbf{(E) }15\sqrt{3}$

 \textbf{Answer: }E



\textbf{Solution 1}

The length of the side opposite to the angle with $120^\circ$ is longest. We denote its value as $x$.

 Because three side lengths form an arithmetic sequence, the middle-valued side length is $\frac{x + 6}{2}$.

 Following from the law of cosines, we have
 \begin{align*} 6^2 + \left( \frac{x + 6}{2} \right)^2  - 2 \cdot 6 \cdot \frac{x + 6}{2} \cdot \cos 120^\circ = x^2 . \end{align*}


 By solving this equation, we get $x = 14$.Thus, $\frac{x + 6}{2} = 10$.

 Therefore, the area of the triangle is
 \begin{align*} \frac{1}{2} \cdot 6 \cdot 10 \cdot \sin 120^\circ = \boxed{\textbf{(E) } 15 \sqrt{3}} . \end{align*}





\textbf{Solution 2}

Let the side lengths of $\triangle ABC$ be $6$, $x$, and $2x-6$, where $6 \le x \le 2x-6$. As $2x-6$ is the longest side, the angle opposite to it will be $120^{\circ}$.

 By the law of Cosine
 \[(2x-6)^2 = 6^2 + x^2 - 2 \cdot 6 \cdot x \cdot \cos 120^{\circ}\]
 \[4x^2 - 24x + 36 = 36 + x^2 + 6x\]
 \[3x^2 - 30x = 0\]
 \[x^2 - 10x = 0\]As $x \neq 0$, $x = 10$.

 Therefore, $[ABC] = \frac{ 6 \cdot 10 \cdot \sin 120^{\circ} }{2} = \boxed{\textbf{(E) } 15 \sqrt{3}}$




\textbf{Solution 3}

Let the side lengths of $\triangle ABC$ be $6$, $6+d$, and $6+2d$, where $6 \le 6+d \le 6+2d$. As $6+2d$ is the longest side, the angle opposite to it will be $120^{\circ}$.

 By the law of Cosine
 \[(6+2d)^2 = 6^2 + (6+d)^2 - 2 \cdot 6 \cdot (6+d) \cdot \cos 120^{\circ}\]
 \[4d^2 + 24d + 36 = 36 + 36 + 12 d + d^2 + 36 + 6d\]
 \[3d^2 + 6d - 72 = 0\]
 \[d^2 + 2d - 24 = 0\]
 \[(d+6)(d-4)=0\]As $d \ge 0$, $d = 4$, $6+d = 10$

 Therefore, $[ABC] = \frac{ 6 \cdot 10 \cdot \sin 120^{\circ} }{2} = \boxed{\textbf{(E) } 15 \sqrt{3}}$




\textbf{Solution 4 (Analytic Geometry)}

Since the triangle's longest side must correspond to the $120^\circ$ angle, the triangle is unique. By analytic geometry, we construct the following plot.

 
                    \begin{figure}[H]
                        \centering
                        \includegraphics[width=0.5\textwidth]{Solution/AMC_12_12B_2023_17_Solution_Figure_1.png}
                    \end{figure}
                    

 We know the coordinates of point $A$ being the origin and $B$ being $(6,0)$. Constructing the line which point $C$ can lay on, here since $\angle B=120^\circ$, $C$ is on the line 
 \[y=\sqrt{3}\left(x-6\right).\]

 I denote $D$ as the perpendicular line from $C$ to $AB$, and assume $CD=k$. Here we know $\triangle BCD$ is a $30^\circ-60^\circ-90^\circ$ triangle. Hence $DC=\sqrt{3}k$ and $BC=2k$.

 Furthermore, due to the arithmetic progression, we know $AC=4k-6$. Hence, in $\triangle ACD$, 
 \[\left(4k-6\right)^{2}=\left(6+k\right)^{2}+3k^{2},\] 
 \[k=5.\]

 Thus, the area is equal to $\frac{1}{2}\cdot 6\cdot \sqrt{3} k=\boxed{\textbf{(E) } 15 \sqrt{3}}$.

 


	\item (\textit{AMC 12 12B 2023}) Last academic year Yolanda and Zelda took different courses that did not necessarily administer the same number of quizzes during each of the two semesters. Yolanda's average on all the quizzes she took during the first semester was $3$ points higher than Zelda's average on all the quizzes she took during the first semester. Yolanda's average on all the quizzes she took during the second semester was $18$ points higher than her average for the first semester and was again $3$ points higher than Zelda's average on all the quizzes Zelda took during her second semester. Which one of the following statements cannot possibly be true?


 $\textbf{(A)}$ Yolanda's quiz average for the academic year was $22$ points higher than Zelda's.


 $\textbf{(B)}$ Zelda's quiz average for the academic year was higher than Yolanda's.


 $\textbf{(C)}$ Yolanda's quiz average for the academic year was $3$ points higher than Zelda's.


 $\textbf{(D)}$ Zelda's quiz average for the academic year equaled Yolanda's.


 $\textbf{(E)}$ If Zelda had scored $3$ points higher on each quiz she took, then she would have had the same average for the academic year as Yolanda.


 
 \textbf{Answer: }A



\textbf{Solution 1}

Denote by $A_i$ the average of person with initial $A$ in semester $i \in \left\{1, 2 \right\}$Thus, $Y_1 = Z_1 + 3$, $Y_2 = Y_1 + 18$, $Y_2 = Z_2 + 3$.

 Denote by $A_{12}$ the average of person with initial $A$ in the full year.Thus, $Y_{12}$ can be any number in $\left( Y_1 , Y_2 \right)$ and $Z_{12}$ can be any number in $\left( Z_1 , Z_2 \right)$.

 Therefore, the impossible solution is$\boxed{\textbf{(A)}~\text{Yolanda's quiz average for the academic year was 22 points higher than Zelda's.}}$




\textbf{Solution 2}

We can use process of elimination by finding possible solutions to answer choices. Let $y_1$ and $y_2$ be the number of quizzes Yolanda took in the first and second semesters, respectively. Define $z_1$ and $z_2$ similarly for Zelda.

 Answer choice B is satisfied by $(y_1,y_2,z_1,z_2) = (289,1,1,289)$.

 Answer choice C and E are both satisfied by $(y_1,y_2,z_1,z_2) = (17,17,17,17)$.

 Answer choice D is satisfied by $(y_1,y_2,z_1,z_2) = (7,5,5,7)$.

 Therefore the answer is $\boxed{\textbf{(A)}~\text{Yolanda's quiz average for the academic year was 22 points higher than Zelda's.}}$




\textbf{Solution 3 (Algebra)}

Let Yolanda's average for semester $1$ be $y_1$, the number of quizzes Yolanda took in semester $1$ be $n_1$, Zelda's average for semester $1$ be $z_1$, the number of quizzes Zelda took in semester $1$ be $k_1$, Yolanda's average for semester $2$ be $y_2$, the number of quizzes Yolanda took in semester $2$ be $n_2$, Zelda's average for semester $2$ be $z_2$, the number of quizzes Zelda took in semester $2$ be $k_2$, Yolanda's average for the entire year be $y$, Zelda's average for the entire year be $z$.

 From the problem we know that

 
 \[y_1 = z_1 + 3, \quad y_2 = z_2 + 3\]

 \[y_2 = y_1 + 18, \quad z_2 = z_1 + 18\]

 \[y = \frac{ y_1 n_1 + y_2 n_2 }{ n_1 + n_2} = \frac{ y_1 n_1 + (y_1 + 18) n_2 }{ n_1 + n_2} = y_1 + \frac{18 n_2 }{ n_1 + n_2}\]

 \[z = \frac{ z_1 k_1 + z_2 k_2 }{ k_1 + k_2} = \frac{ z_1 k_1 + (z_1 + 18) k_2 }{ k_1 + k_2} = z_1 + \frac{ 18 k_2 }{ k_1 + k_2}\]

 \[y - z = y_1 + \frac{18 n_2 }{ n_1 + n_2} - z_1  - \frac{ 18 k_2 }{ k_1 + k_2} = y_1 - z_1 + 18 \left( \frac{n_2 }{ n_1 + n_2} - \frac{k_2 }{ k_1 + k_2} \right) = 3 + 18 \left( \frac{n_2 }{ n_1 + n_2} - \frac{k_2 }{ k_1 + k_2} \right)\]
$\frac{n_2 }{ n_1 + n_2}$ at most is $1$, $\frac{k_2 }{ k_1 + k_2}$ is at least $0$, meaning that $\frac{n_2 }{ n_1 + n_2} - \frac{k_2 }{ k_1 + k_2}$ is at most $1$.

 Therefore, 
 \[y - z \le 3 + 18 = 21\]

 Hence, $\boxed{\textbf{(A)}~\text{Yolanda's quiz average for the academic year was 22 points higher than Zelda's}}$ is not possible.

 



\textbf{Solution 4}

Denote $y_1$ and $y_2$ as the quiz averages of Yolanda in the $1$st and $2$nd semesters, respectively. Similarly, denote $z_1$ and $z_2$ as the quiz averages of Zelda in the $1$st and $2$nd semesters. 

 We have $y_1 = z_1 + 3$, so $y_1 - 3 = z_1$. We also know that $y_2 = 18 +y_1 = 3 + z_2$, implying $z_2 = 15 + y_1$.

 The average quiz scores for both students must lie between the averages of each semester, i.e $\mathrm{Avg_{y}}\in[y_1, 18+y_1],$ and $\mathrm{Avg_{z}}\in[y_1-3, y_1+15].$ Since $z_2>z_1$, $y_2 > y_1$, and $y_1>z_1$, we have $\min\{\mathrm{Avg}\} = \min\{\mathrm{Avg}_z\} = z_1$ and $\max\{\mathrm{Avg}\} = \max\{\mathrm{Avg}_y\} = y_2$. Therefore the maximum difference between the two yearly averages is 

 
 \[|y_2 - z_1| = |(y_1 + 18)-(y_1-3)| = 21 < 22.\]

 Therefore, $\boxed{\textbf{(A)}~\text{Yolanda's quiz average for the academic year was 22 points higher than Zelda's}}$ is not possible.

 -Benedict T (countmath1)

 



\textbf{Solution 5}

Let\begin{align*}\text{Yolanda's 1st Semester Average} &= \frac{a_1+\dots+a_n}{n} \\\text{Yolanda's 2nd Semester Average} &= \frac{a_1'+\dots+a_n'}{n'} \\\text{Zelda's 1st Semester Average} &= \frac{b_1+\dots+b_m}{m} \\\text{Zelda's 2nd Semester Average} &= \frac{b_1'+\dots+b_m'}{m'}.\end{align*}


 According to the problem, $\frac{a_1+\dots+a_n}{n}+18=\frac{a_1'+\dots+a_n'}{n'}$.

 Moreover, $\frac{a_1+\dots+a_n}{n}=\frac{b_1+\dots+b_m}{m}+3$ and $\frac{a_1'+\dots+a_n'}{n'}=\frac{b_1'+\dots+b_m'}{m'}+3$ are true.

 In another words,
 \[\frac{a_1+\dots+a_n}{n}=\frac{a_1'+\dots+a_n'-18n'}{n'}=\frac{b_1+\dots+b_m+3m}{m}= \frac{b_1'+\dots+b_m'-15m'}{m'}.\]

 The relationship leads to
 \[\frac{a_1+\dots+a_n+a_1'+\dots+a_n'-18n'}{n+n'}=\frac{b_1+\dots+b_m+b_1'+\dots+b_m'+3m-15m'}{m+m'}, \text{or}\]
 \[\frac{a_1+\dots+a_n+a_1'+\dots+a_n'}{n+n'}-\frac{18n'}{n'+n}-\frac{3m}{m+m'}+\frac{15m'}{m+m'}=\frac{b_1+\dots+b_m+b_1'+\dots+b_m'}{m+m'}.\]For choice $(A)$ to be true, $-\frac{18n'}{n'+n}-\frac{3m}{m+m'}+\frac{15m'}{m+m'}$ must be $-22$.

 
Let $n'=kn$ and $m'=k'm$.
 \[-\frac{18n'}{n'+n}-\frac{3m}{m+m'}+\frac{15m'}{m+m'}=-\frac{18kn}{n(k+1)}-\frac{3m}{m(k'+1)}+\frac{15k'm}{m(k'+1)}=-\frac{18kn}{n(k+1)}+\frac{3m(5k'-1)}{m(k'+1)}\]$n,n',m,m',k\text{ and }k'$ must be non-zero integer to satisfy the condition in the question. While the equation above is a multiple of three, 22 is not. Therefore, choice $\boxed{(A) \text{ cannot be true.}}$



	\item (\textit{AMC 12 12B 2023}) Each of $2023$ balls is placed in one of $3$ bins. Which of the following is closest to the probability that each of the bins will contain an odd number of balls?


 $\textbf{(A) } \frac{2}{3} \qquad \textbf{(B) } \frac{3}{10} \qquad \textbf{(C) } \frac{1}{2} \qquad \textbf{(D) } \frac{1}{3} \qquad \textbf{(E) } \frac{1}{4}$

 \textbf{Answer: }E



\textbf{Solution 1 (Systematic Algebraic Way)}

Suppose the numbers are $a_1$, $a_2$, and $a_3$. First, we try to calculate the amount of ways for all three balls to be placed in a bin so the number of balls in each bin is odd. $a_1+a_2+a_3=2023$ and each bin has at least one ball because they are positive odd numbers. Changing the equation, we see that $\frac{[(a_1)+1]}{2}+\frac{[(a_2)+1]}{2}+\frac{[(a_3)+1]}{2}=\frac{(2023+3)}{2}$. Let $a_1+1=2b_1$, $a_2+1=2b_2$, and $a_3+1=2b_3$. Thus $b_1+b_2+b_3=1013$. We can also see that $b_1$, $b_2$, and $b_3$ are all positive. Using the positive version of stars and bars, we get ${1013-1 \choose 3-1}$ = ${1012 \choose 2}$ choices. 

 Now, we want to find the total amount of cases. Using the non-negative version of stars and bars, we find that the total is ${2023+3-1 \choose 3-1}$=${2025 \choose 2}$. 

 Now we need to calculate ${1012 \choose 2}$/${2025 \choose 2}$, which is just $\frac{1012 \cdot 1011 \cdot 2}{2 \cdot 2025 \cdot 2024}$. Cancelling the twos, we get $\frac{1012 \cdot 1011}{2025 \cdot 2024}$. This is roughly equal to $\frac{1}{4}$. The answer is $\boxed{\textbf{(E)} \frac{1}{4}}$.

 



\textbf{Solution 2 (Pairing)}

There are \( \left\lceil\frac{2023}{2}\right\rceil = 1012 \) odd numbers between 1 and 2023.

 Out of these, there are \( \left\lfloor\frac{1012}{3}\right\rfloor = 337 \) ways to evenly distribute each ball into the first bin, then, WLOG, we say there is \( 337 - 2 = 335 \) for the second and \( 337 + 2 = 339 \) for the third bin.

 Out of these, there are \( \left\lfloor\frac{337}{2}\right\rfloor = 168 \) odd sums of pairs for the first bin, then \( \left\lfloor\frac{335}{2}\right\rfloor = 167 \) for the second, and \( \left\lfloor\frac{339}{2}\right\rfloor = 169 \) for the third. The sum of all of these is \( 168 + 167 + 169 = 504 \) total distributions.

 From here, we need to find the probability over 2023 total distributed balls, which is simply \( 504/2023 \approx 500/2000 = 1/4 \Rightarrow \) $\boxed{\textbf{(E)} \frac{1}{4}}$.

 This is not an exact probability. It is simply a solution for general competition math. For the exact way to calculate the probability (highly not recommended), please refer to Solution 8.

 



\textbf{Solution 3}

We first examine the possible arrangements for parity of number of balls in each box for $2022$ balls.

 If a $E$ denotes an even number and a $O$ denotes an odd number, then the distribution of balls for $2022$ balls could be $EEE,EOO,OEO,$ or $OOE$. With the insanely overpowered magic of cheese, we assume that each case is about equally likely.

 From $EEE$, it is not possible to get to all odd by adding one ball; we could either get $OEE,EOE,$ or $EEO$. For the other $3$ cases, though, if we add a ball to the exact right place, then it'll work.

 For each of the working cases, we have $1$ possible slot the ball can go into (for $OEO$, for example, the new ball must go in the center slot to make $OOO$) out of the $3$ slots, so there's a $\dfrac13$ chance. We have a $\dfrac34$ chance of getting one of these working cases, so our answer is $\dfrac34\cdot\dfrac13=\boxed{\textbf{(E) }\dfrac14.}$




\textbf{Solution 4}

2023 is an arbitrary large number. So, we proceed assuming that an arbitrarily large number of balls have been placed.

 For an odd-numbered amount of balls case, the 3 bins can only be one of these 2 combinations:

 $OEE$ ($OEE$,$EOE$,$EEO$)

 $OOO$ ($OOO$)

 Let the probability of achieving the $OOO$ case to be $P(OOO) = p$ and any of the $OEE$ permutations to be $P(OEE) = 1-p$.

 Because the amount of balls is arbitrarily large, $P(OOO) = p$ even after another two balls are be placed.

 There are two cases for which placing another two balls results in $OOO$:

 $OOO$: The two balls are placed in the same bin ($OOO\to OOE\to OOO$)

 $OEE$: The two balls are placed in the two even bins ($OEE\to OOE \to OOO$)

 So, 

 $P(OOO) = P(OOO) * \frac{1}{3} + P(OEE) * \frac{2}{3} * \frac{1}{3}$
$p = p * \frac{1}{3} + (1-p) * \frac{2}{3} * \frac{1}{3}$
$\frac{8}{9}p = \frac{2}{9}$
$p = \boxed {\textbf {(E)} \frac {1}{4}}$




\textbf{Solution 5}

We use the generating functions approach to solve this problem.Define $\Delta = \left\{ \left( a, b, c \right) \in \Bbb Z_+: a+b+c = 2023 \right\}$.

 We have
 \[ \left( x + y + z \right)^{2023} = \sum_{(a,b,c) \in \Delta} \binom{2023}{a,b,c} x^a y^b z^c . \]

 First, we set $x \leftarrow 1$, $y \leftarrow 1$, $z \leftarrow 1$.We get
 \[ 3^{2023} = \sum_{(a,b,c) \in \Delta} \binom{2023}{a,b,c} 1 . \hspace{1cm} (1) \]

 Second, we set $x \leftarrow 1$, $y \leftarrow -1$, $z \leftarrow 1$.We get
 \[ 1 = \sum_{(a,b,c) \in \Delta} \binom{2023}{a,b,c} (-1)^b . \hspace{1cm} (2) \]

 Third, we set $x \leftarrow 1$, $y \leftarrow 1$, $z \leftarrow -1$.We get
 \[ 1 = \sum_{(a,b,c) \in \Delta} \binom{2023}{a,b,c} (-1)^c . \hspace{1cm} (3) \]

 Fourth, we set $x \leftarrow 1$, $y \leftarrow -1$, $z \leftarrow -1$.We get
 \[ -1 = \sum_{(a,b,c) \in \Delta} \binom{2023}{a,b,c} (-1)^{b+c} . \hspace{1cm} (4) \]

 Taking $\frac{(1)-(2) - (3)+(4)}{4}$, we get
 \begin{align*} \frac{3^{2023} - 1 - 1 + (-1)}{4} & = \frac{1}{4} \sum_{(a,b,c) \in \Delta} \binom{2023}{a,b,c}  \left( 1 - (-1)^b - (-1)^c + (-1)^{b+c} \right) \\ & = \frac{1}{4} \sum_{(a,b,c) \in \Delta} \binom{2023}{a,b,c}  \left( 1 - (-1)^b \right) \left( 1 - (-1)^c \right) \\ & = \sum_{\substack{(a,b,c) \in \Delta \\ a, b, c \mbox{ are odds}}} \binom{2023}{a,b,c} . \end{align*}


 The last expression above is the number of ways to get all three bins with odd numbers of balls.Therefore, this happens with probability
 \begin{align*} \frac{\frac{3^{2023} - 1 - 1 + (-1)}{4}}{3^{2023}} & \approx \boxed{\textbf{(E) } \frac{1}{4}}. \end{align*}





\textbf{Solution 6}

Four even-odd splittings divides $2023$ in to three, namely $(O,O,O)$, $(E,E,O)$, $(E,O,E)$, and $(O,E,E)$. Here if we define a "move" as relocated one ball, then we will notice in each case, that a random "move" will be evenly likely to be one of the other three splittings. Hence by Group theory (or by intuition), we will find the structure of this splitting is $V_4$ group, and it's symmetric for all four elements in this Group.

 Thus, no matter what is the initial starting point, four cases will be evenly likely to appear when repeated many times. The answer is $\boxed{\textbf{(E)} \frac{1}{4}}$.

 



\textbf{Solution 7 (Exact Probability)}

Let $p(n)$ be the probability of all 3 bins having odd number of balls when we randomly put $2n+1$ balls into these bins. We need find out $p(1011)$. Clearly $p(0)=0$. For $n>1$, there are only two scenarios for the first $2n-1$ balls: \textbf{Case 1} all 3 bins are odd; \textbf{Case 2} only 1 bin is odd and the other 2 bins are even.

 For case 1 which has a probability $p(n-1)$, in order to get all 3 bins to be odd, we need to put the last 2 balls in the same bin. This has a probability of $\frac{1}{3}$. 

 For case 2 which has a probability of $1-p(n-1)$, in order to get all 3 bins to be odd, we need to put the last 2 balls into the 2 even bins. The first one has a probability of $\frac{2}{3}$ to be put in an even bin, and the second has a probability of $\frac{1}{3}$ to be put the in the last even bin, so the probability is $\frac{2}{9}$ to get all 3 bins to be odd in this case.

 So we have $p(n)=\frac{1}{3} p(n-1) + \frac{2}{9} (1-p(n-1)) = \frac{1}{9} p(n-1) + \frac{2}{9}$.

 Solving for the fixed point: $x=\frac{1}{9} x + \frac{2}{9}$ yields $x=\frac{1}{4}$, and this helps us the rewrite the above as $p(n) - \frac{1}{4} = \frac{1}{9} (p(n-1)  - \frac{1}{4} )$.

 Therefore $p(n)-\frac{1}{4} = \frac{1}{9^n} (p(0) - \frac{1}{4})$. Using $p(0)=0$, we have $p(n)=\frac{1}{4} (1 - \frac{1}{9^n})$.

 The probability for 2023 balls is $\frac{1}{4} (1 - \frac{1}{9^{1011}})$, and this is closest to $\boxed{\textbf{(E)} \frac{1}{4}}$.

 



\textbf{Solution 8 (Recursion, rigorous)}

We can use recursion to figure out an expression that determines the probability of all the bins being odd.

 Define $P_n$ to be the probability that with $2n-1$ balls, randomly inserting them makes all 3 bins have an odd count.We would first want to find $P_{n+1}=f(P_n)$ for some function $f$.

 Now, we can say that if we have an O-O-O combination, where O denotes odd and E denotes even, then to add 2 more balls to keep it O-O-O, we need to put them in the same bin. This would result in a $\frac{1}{3}$ probability of staying O-O-O.

 If we have an E-E-O combination, then we need to add 2 more balls, one in each of the even bins. This would result in a $\frac{2}{9}$ probability of turning to O-O-O.Now, we can write a function of $P_{n+1}$ in terms of $P_n$.

 $P_{n+1}=\frac{1}{3}P_n+\frac{2}{9}-\frac{2}{9}P_n=\frac{1}{9}P_n+\frac{2}{9}$

 Now, we want to express $P_{n+1}=g(n)$ for some function $g$. Here, the $\frac{2}{9}$ is a problem because now we can't just put it as a geometric sequence. However, we can rewrite it as followed:

 $P_{n+1}=\frac{1}{9}P_n+\frac{2}{9}$
$P_{n+1}-\frac{1}{4}=\frac{1}{9}(P_n-\frac{1}{4})$

 
Now, we can let $Q_n=P_n-\frac{1}{4}$. We can rewrite that as $Q_{n+1}=\frac{1}{9}Q_n$.

 We can note that $Q_1=P_1-\frac{1}{4}$. Since $P_1=0$ (you only have one ball, so two of the bins have to be even), $Q_1=-\frac{1}{4}$. Therefore, $Q_n=Q_1(\frac{1}{9})^{n-1}=-\frac{1}{4} \cdot (\frac{1}{9})^{n-1}$.

 Substituting, we can find that $P_n=-\frac{1}{4} \cdot (\frac{1}{9})^{n-1}+\frac{1}{4}$.

 2023 is the 1012th odd number, so we can substitute $n=1012$ into the equation.

 $-\frac{1}{4} \cdot (\frac{1}{9})^{1011}+\frac{1}{4} \approx 0+\frac{1}{4}=\frac{1}{4}$.

 Therefore, the answer is $\boxed{\textbf{(E) }\frac{1}{4}}$




\textbf{Solution 9 (Generating Functions instead of stars and bars)}

Define 
 \begin{align*} [x^{2023}]G(x)&=(x^1 + x^3 + ... + x^{2023} + ...)^3 \\ [x^{2023}]G(x)&= x^3(1+x^2+x^4+...+x^{2022}+...)^3 \\ [x^{2020}]G(x)&= (\frac{1}{1-x^2})^3 \\ [x^{2020}]G(x)&= (1-x^2)^{-3} \\ [x^{2020}]G(x)&= (1 + \frac{3 \cdot 2}{2 \cdot 1}(x^2)^1 + ... \frac{1012 \cdot 1011}{2 \cdot 1}(x^2)^{1010}) \\ [x^{2020}] &= \frac{1012 \cdot 1011}{2 \cdot 1} \end{align*}


 Similarly we can compute the total number of cases which is: $\frac{2025\cdot 2022}{2 \cdot 1}$

 So, the probability is: 
 \begin{align*} &= \frac {\frac{1012 \cdot 1011}{2 \cdot 1}}{\frac{2025\cdot 2024}{2 \cdot 1}} \\ &= \frac{1012 \cdot 1011}{2025 \cdot 2024} \\ &\approx \frac{1}{4} \\ \end{align*}


 Therefore, the answer is $\boxed{\textbf{(E) }\frac{1}{4}}$




\textbf{Solution 10 (Short On Time)}

One must not forget that the balls are placed randomly. With intuition, we can assume that the probability of getting $OOO$, $OOE$, $OEE$, and $EEE$, where $O$ denotes odd numbers and $E$ denotes even numbers, would nearly be the same. Since there are $4$ possible outcomes and $1$ of them is desired, we can conclude that the likelihood of getting $OOO$ is $\frac{1}{4}$ and that the correct answer choice is $\boxed{\textbf{(E) }\frac{1}{4}}$.

 Note: Use this type of reasoning only when a logical solution doesn't work out for you.

 Another Note (not by Aqf243): This may not be the best solution because 2023 is an odd number. Hence, realistically, the probability of getting $OOE$ is 0, allowing us to remove it from our list. This is because adding 2 odd numbers (represented as 2 numbers $x \in \{ x \mid x \bmod 2 = 1 \}$) always results in an even number.

 
$\because \newline$ $x_1=1\mod{2}$
$x_2=1\mod{2}$
$\therefore \\ x_1+x_2=1+1 \mod{2}=2 \mod{2}= 0 \mod{2}$

 for any number $x\in{{}x\mid x\mod{2}=0{}}$, $x$ is an even number.

 Similarly, adding two even numbers creates another even number

 $x_1=0 \mod{2} \\ x_2=0 \mod{2}$$x_1+x_2=0+0 \mod{2}=0 \mod{2}$

 So adding 2 odd numbers an an even number ALWAYS results in an even number. 

 Basically, Aqf243 got a little lucky in this situation...

 



\textbf{Solution 11 (Notice a pattern)}

Let the bins be called $a$, $b$, and $c$. We want the number of odd positive integer solutions that satisfy $a + b + c = 2023$. We know from stars and bars that the total number of integer solutions is $\binom{2023 + 3 - 1}{3 - 1} = \binom{2025}{2}$. If $a = 1$, $b$ can be any odd number from $1$ to $2021$. There are $(2021 + 1)/2 = 1011$ choices for $b$. If $a = 3$, there would be $1010$ choices for $b$, and so on. $c$ is always constrained by $a$ and $b$ so our answer is just the number of choices $b$ has for each $a$, which is the sum of all integers from $1$ to $1011 = \frac{1011 \cdot 1012}{2}$. So we get $\frac{1011 \cdot 1012/2}{\binom{2025}{2}}$ (around $0.249$), which is closest to $0.25$ or $\boxed{\textbf{(E) }\frac{1}{4}}$.

 \textbf{NOTE:} The probability we came to in this solution and in other solutions that use Stars and Bars is not exact. Stars and Bars treats all outcomes to be equally likely when in reality some distributions are more likely than others (see clarification above), but the reason we still use it here is because it gives us a good average and the answer choices are pretty far apart so it’s not very significant to our approximation.

 


	\item (\textit{AMC 12 12B 2023}) Cyrus the frog jumps $2$ units in a direction, then $2$ more in another direction. What is the probability that he lands less than $1$ unit away from his starting position?


 $\textbf{(A)}~\frac{1}{6}\qquad\textbf{(B)}~\frac{1}{5}\qquad\textbf{(C)}~\frac{\sqrt{3}}{8}\qquad\textbf{(D)}~\frac{\arctan \frac{1}{2}}{\pi}\qquad\textbf{(E)}~\frac{2\arcsin \frac{1}{4}}{\pi}$

 \textbf{Answer: }E



\textbf{Solution 1}


                    \begin{figure}[H]
                        \centering
                        \includegraphics[width=0.5\textwidth]{Solution/AMC_12_12B_2023_20_Solution_Figure_1.png}
                    \end{figure}
                    Let Cyrus's starting position be $S$. WLOG, let the place Cyrus lands at for his first jump be $O$. From $O$, Cyrus can reach all the points on $\odot O$. The probability that Cyrus will land less than $1$ unit away from $S$ is $\frac{4 \alpha }{ 2 \pi}$.

 
 \[\sin \alpha = \frac{ \frac12 }{2} = \frac14, \quad \alpha = \arcsin \frac14\]

 Therefore, the answer is

 
 \[\frac{4 \arcsin \frac14 }{ 2 \pi} = \boxed{\textbf{(E) } \frac{2 \arcsin \frac{1}{4}}{\pi}}\]




\textbf{Solution 2}

Denote by $A_i$ the position after the $i$th jump.Thus, to fall into the region centered at $A_0$ and with radius 1, $\angle A_2 A_1 A_0 < 2 \arcsin \frac{1/2}{2} = 2 \arcsin \frac{1}{4}$.

 Therefore, the probability is
 \[ \frac{2 \cdot 2 \arcsin \frac{1}{4}}{2 \pi} = \boxed{\textbf{(E) } \frac{2 \arcsin \frac{1}{4}}{\pi}}. \]




\textbf{Solution 3(coord bash)}

Let the orgin be the starting point of frog. Then, WLOG assume that after the first jump, it is at the point (2,0). Then, the range of all possible places the frog can jump to at its second jump is the circle with equation $(x-2)^2+y^2=2^2$.If it landed $1$ unit within its starting point (the orgin), then it is inside the circle $x^2+y^2=1$. We clearly want the intersection point. So we're trying to solve the system of equations $x^2+y^2=1$ and $(x-2)^2+y^2=2^2$. We have $x=\frac{1}{4}$, so $y=\pm\frac{\sqrt{15}}{4}$. Therefore, our desired answer would be $\frac{\arcsin{\frac{\sqrt{15}}{8}}}{\pi}$ (the angle we want divided by $2\pi$). Since

 
 \[\arcsin{\frac{\sqrt{15}}{8}}=\arccos{\frac{7}{8}}=\arccos{(1-2 \cdot (\frac{1}{4})^{2})}=2\arcsin{\frac{1}{4}}\]

 where the last step holds by the double angle formula, our answer is $\boxed{(E)\frac{2\arcsin{\frac{1}{4}}}{\pi}}$.~\href{https://artofproblemsolving.com/wiki/index.php?title=Ddk001}{Ddk001}

 



\textbf{Solution 4 - Law of Cosines and Double Angle Formula}

Let $A$ be Cyrus's starting point, $B$ be the first point he jumps to ($AB = 2$), and $C$ be the second point he jumps to ($BC = 2$). Let angle $ABC$ be $k$, such that $AC = 1$. The probability of $AC < 1$ would therefore be $\frac{2k}{360}$ (since $C$ could be on either side of $AB$ so there are two possible areas of having $AC < 1$) which simplifies to $\frac{k}{180}$. Converting to radians gives us $\frac{k}{\pi}$. To find $k$, we use the law of cosines.

 
 \[AC^2 = AB^2 + BC^2 - 2 \cdot AB \cdot BC \cdot \cos k\] 
 \[1^2 = 2^2 + 2^2 - 2 \cdot 2 \cdot 2 \cdot \cos k\] 
 \[1 = 4 + 4 - 8 \cos k\] 
 \[8 \cos k = 7\] 
 \[\cos k = \frac{7}{8}\] 
 \[k = \arccos\left(\frac{7}{8}\right) = \arcsin\left(\sqrt{1 - \left(\frac{7}{8}\right)^2}\right) = \arcsin\left(\sqrt{\frac{15}{64}}\right) = \arcsin\left(\frac{1}{4} \sqrt{\frac{15}{4}}\right) = \arcsin\left(\frac{1}{4} \sqrt{1 - \left(\frac{1}{4}\right)^2}\right)\] 
 \[= \arcsin\left(\sin\left(\arcsin\left(\frac{1}{4}\right)\right) \cos\left(\arcsin\left(\frac{1}{4}\right)\right)\right) = \arcsin\left(\sin\left(2 \arcsin\left(\frac{1}{4}\right)\right)\right) = 2 \arcsin\left(\frac{1}{4}\right)\]

 The probability is 
 \[\frac{k}{\pi} = \frac{2 \arcsin\left(\frac{1}{4}\right)}{\pi}\]

 which is $E$.

 



\textbf{Solution 5 - Complex Number and Double Angle Formula}


                    \begin{figure}[H]
                        \centering
                        \includegraphics[width=0.5\textwidth]{Solution/AMC_12_12B_2023_20_Solution_Figure_2.png}
                    \end{figure}
                     The frog jumps 2 units in one direction, which we can represent as $z_1 = 2e^{ia}$, where $a$ is the angle in radians. Then, the frog jumps another 2 units in a different direction, represented as $z_2 = 2e^{ib}$, where $b$ is the angle in radians.

 The total displacement from the starting position is $z_3 = z_1 + z_2 = 2e^{ia} + 2e^{ib}$.

 
We need to find the condition under which $|z_3| <= 1$. This translates to:
 \[|2e^{ia} + 2e^{ib}| <= 1\]

 \[|e^{ia} + e^{ib}| <= \frac{1}{2}\]

 \[e^{ia} + e^{ib} = \cos(a) + i\sin(a) + \cos(b) + i\sin(b)\]
 \[= (\cos(a) + \cos(b)) + i(\sin(a) + \sin(b))\]

 The magnitude squared of this sum is:
 \[|e^{ia} + e^{ib}|^2 = (\cos(a) + \cos(b))^2 + (\sin(a) + \sin(b))^2\]
 \[= \cos^2(a) + 2\cos(a)\cos(b) + \cos^2(b) + \sin^2(a) + 2\sin(a)\sin(b) + \sin^2(b)\]
 \[= 2 + 2(\cos(a)\cos(b) + \sin(a)\sin(b))\]
 \[= 2 + 2\cos(a - b)\]

 Thus:
 \[2 + 2\cos(a - b) <= \frac{1}{4}\]
 \[2\cos(a - b) <= -\frac{7}{4}\]
 \[\cos(a - b) <= -\frac{7}{8}\]

 Therefore, our desired answer would be $\frac{\arccos{\frac{7}{8}}}{\pi}$ (the angle we want divided by $2\pi$). Since

 
 \[\arccos{\frac{7}{8}}=\arccos{(1-2 \cdot (\frac{1}{4})^{2})}=2\arcsin{\frac{1}{4}}\]

 where the last step holds by the double angle formula, our answer is $\boxed{(E)\frac{2\arcsin{\frac{1}{4}}}{\pi}}$.

 


	\item (\textit{AMC 12 12B 2023}) A lampshade is made in the form of the lateral surface of the frustum of a right circular cone. The height of the frustum is $3\sqrt3$ inches, its top diameter is $6$ inches, and its bottom diameter is $12$ inches. A bug is at the bottom of the lampshade and there is a glob of honey on the top edge of the lampshade at the spot farthest from the bug. The bug wants to crawl to the honey, but it must stay on the surface of the lampshade. What is the length in inches of its shortest path to the honey?


 $\textbf{(A) } 6 + 3\pi\qquad \textbf{(B) }6 + 6\pi\qquad \textbf{(C) } 6\sqrt3 \qquad \textbf{(D) } 6\sqrt5 \qquad \textbf{(E) } 6\sqrt3 + \pi$

 \textbf{Answer: }E



\textbf{Solution}

We augment the frustum to a circular cone.Denote by $O$ the apex of the cone.Denote by $A$ the bug and $B$ the honey.

 By using the numbers given in this problem, the height of the cone is $6 \sqrt{3}$.Thus, $OA = 12$ and $OB = 6$.

 We unfold the lateral face. So we get a circular sector.The radius is 12 and the length of the arc is $12\pi$.Thus, the central angle of this circular sector is $180^\circ$.

 Because $A$ and $B$ are opposite in the original frustum, in the unfolded circular cone, $\angle AOB = \frac{180^\circ}{2} = 90^\circ$.

 Notice that a feasible path between $A$ and $B$ can only fall into the region with the range of radii between $OB = 6$ and $OA = 12$.Therefore, we cannot directly connect $A$ and $B$ and must make a detour.Denote by $AC$ a tangent to the circular sector with radius 6 that meets it at point $C$.Therefore, the shortest path between $A$ and $B$ consists of a segment $AC$ and an arc from $C$ to $B$.

 Because $OA = 12$, $OC = 6$ and $\angle OCA = 90^\circ$, we have $AC = \sqrt{OA^2 - OC^2} = 6 \sqrt{3}$ and $\angle AOC = 60^\circ$.This implies $\angle COB = \angle AOB - \angle AOC = 30^\circ$.Therefore, the length of the arc between $C$ and $B$ is $OB \cdot \pi \cdot  \frac{\angle COB}{180^\circ} = \pi$.Therefore, the shortest distance between $A$ and $B$ is $\boxed{\textbf{(E) } 6 \sqrt{3} + \pi}$.

 


	\item (\textit{AMC 12 12B 2023}) A real-valued function $f$ has the property that for all real numbers $a$ and $b,$
 \[f(a + b)  + f(a - b) = 2f(a) f(b).\]Which one of the following cannot be the value of $f(1)?$
$\textbf{(A) } 0 \qquad \textbf{(B) } 1 \qquad \textbf{(C) } -1 \qquad \textbf{(D) } 2 \qquad \textbf{(E) } -2$

 \textbf{Answer: }E



\textbf{Solution 1}

Substituting $a = b$ we get 
 \[f(2a) + f(0) = 2f(a)^2\]Substituting $a= 0$ we find 
 \[2f(0) = 2f(0)^2 \implies f(0) \in \{0, 1\}.\]This gives 
 \[f(2a) = 2f(a)^2 - f(0) \geq 0-1\]Plugging in $a = \frac{1}{2}$ implies $f(1) \geq -1$, so answer choice $\boxed{\textbf{(E) -2}}$ is impossible. 

 



\textbf{Solution 2}

First, we set $a \leftarrow 0$ and $b \leftarrow 0$.Thus, the equation given in the problem becomes $[ f(0) + f(0) = 2 f(0) \times f(0) . ]$

 Thus, $f(0) = 0$ or 1.

 Case 1: $f(0) = 0$.

 We set $b \leftarrow 0$.Thus, the equation given in the problem becomes$[ 2 f(a) = 0 . ]$

 Thus, $f(a) = 0$ for all $a$.

 Case 2: $f(0) = 1$.

 We set $b \leftarrow a$.Thus, the equation given in the problem becomes
 \[[ f(2a) + 1 = 2 \left( f(a) \right)^2. ]\]

 Thus, for any $a$,
 \begin{align*} f(2a) & = -1 + 2 \left( f(a) \right)^2 \\ & \geq -1 . \end{align*}


 Therefore, an infeasible value of $f(1)$ is$\boxed{\textbf{(E) -2}}.$




\textbf{Solution 3}

The relationship looks suspiciously like a product-to-sum identity. In fact,
 \[\cos(\alpha)\cos(\beta) = \frac{1}{2}(\cos(\alpha-\beta)+\cos(\alpha+\beta))\]which is basically the relation. So we know that $f(x) = \cos(x)$ is a valid solution to the function. However, if we define $x=ay,$ where $a$ is arbitrary, the above relation should still hold for $f(x) = \cos(ay) = \cos(a(1))$ so any value in $[-1,1]$ can be reached, so choices $A,B,$ and $C$ are incorrect.

 In addition, using the similar formula for hyperbolic cosine, we know
 \[\cosh(\alpha)\cosh(\beta) = \frac{1}{2}(\cosh(\alpha-\beta)+\cosh(\alpha+\beta))\]The range of $\cosh(ay)$ is $[1,\infty)$ so choice $D$ is incorrect.

 Therefore, the remaining answer is choice $\boxed{\textbf{(E) -2}}.$




\textbf{Solution 4 (Recognition)}

Similar to solution 3 we know that $\cos (a - b) \cos (a + b) = \cos (a)\cos(b).$ Doing testing gives us that any function in the form $\cos(bx)$ works. We now see that only 2 and -2 are the candidates to be impossible. Solving the functional equation the regular way, we see that $f(0) = 1, f(-x) = x$. We eventually get a squared term (from even more testing), and see that $\boxed{-2 (E)}$ is the answer because negative square roots are impossible.

 


	\item (\textit{AMC 12 12B 2023}) When $n$ standard six-sided dice are rolled, the product of the numbers rolled can be any of $936$ possible values. What is $n$?


 $\textbf{(A)}~11\qquad\textbf{(B)}~6\qquad\textbf{(C)}~8\qquad\textbf{(D)}~10\qquad\textbf{(E)}~9$

 \textbf{Answer: }A



\textbf{Solution 1}

We start by trying to prove a function of $n$, and then we can apply the function and equate it to $936$ to find the value of $n$.

 It is helpful to think of this problem in the format $(1+2+3+4+5+6) \cdot (1+2+3+4+5+6) \dots$. Note that if we represent the scenario in this manner, we can think of picking a $1$ for one factor and then a $5$ for another factor to form their product - this is similar thinking to when we have the factorized form of a polynomial. Unfortunately this is not quite accurate to the problem because we can reach the same product in many ways: for example for $n=2$, $4$ can be reached by picking $1$ and $4$ or $2$ and $2$. However, this form gives us insights that will be useful later in the problem.

 Note that there are only $3$ primes in the set $\{1,2,3,4,5,6\}$: $2,3,$ and $5$. Thus if we're forming the product of possible values of a dice roll, the product has to be written in the form $2^h \cdot 3^i \cdot 5^j$ (the choice of variables will become clear later), for integer nonnegative values $h,i,j$. So now the problem boils down to how many distinct triplets $(h,i,j)$ can be formed by taking the product of $n$ dice values. 

 We start our work on representing $j$: the powers of $5$, because it is the simplest in this scenario because there is only one factor of $5$ in the set. Because of this, having $j$ fives in our prime factorization of the product is equivalent to picking $j$ factors from the polynomial $(1+\dots + 6) \cdots$ and choosing each factor to be a $5$. Now that we've selected $j$ factors, there are $n-j$ factors remaining to choose our powers of $3$ and $2$. 

 Suppose our prime factorization of this product contains $i$ powers of $3$. These powers of $3$ can either come from a $3$ factor or a $6$ factor, but since both $3$ and $6$ contain only one power of $3$, this means that a product with $i$ powers of $3$ corresponds directly to picking $i$ factors from the polynomial, each of which is either $3$ or $6$ (but this distinction doesn't matter when we consider only the powers of $3$. 

 Now we can reframe the problem again. Our method will be as follows: Suppose we choose an arbitrary pair $(i,j)$ that match the requirements, corresponding to the number of $3$'s and the number of $5$'s our product will have. Then how many different $h$ values for the powers of $2$ are possible?

 In the $i+j$ factors we have already chosen, we obviously can't have any factors of $2$ in the $j$ factors with $5$. However, we can have a factor of $2$ pairing with factors of $3$, if we choose a $6$. The maximal possible power of $2$ in these $i$ factors is thus $2^i$, which occurs when we pick every factor to be $6$. 

 We now have $n-i-j$ factors remaining, and we want to allocate these to solely powers of $2$. For each of these factors, we can choose either a $1,2,$ or $4$. Therefore the maximal power of $2$ achieved in these factors is when we pick $4$ for all of them, which is equivalent to $2^{2\cdot (n-i-j)}$.

 Now if we multiply this across the total $n$ factors (or $n$ dice) we have a total of $2^{2n-2i-2j} \cdot 2^i = 2^{2n-i-2j}$, which is the maximal power of $2$ attainable in the product for a pair $(i,j)$. Now note that every power of $2$ below this power is attainable: we can simply just take away a power of $2$ from an existing factor by dividing by $2$. Therefore the powers of $2$, and thus the $h$ value ranges from $h=0$ to $h=2n-i-2j$, so there are a total of $2n+1-i-2j$ distinct values for $h$ for a given pair $(i,j)$. 

 Now to find the total number of distinct triplets, we must sum this across all possible $i$s and $j$s. Lets take note of our restrictions on $i,j$: the only restriction is that $i+j \leq n$, since we're picking factors from $n$ dice.

 
 \[\sum_{i+j\leq n}^{} 2n+1-i-2j = \sum_{i+j \leq n}^{} 2n+1 - \sum_{i+j \leq n}^{} i+2j\]

 We start by calculating the first term. $2n+1$ is constant, so we just need to find out how many pairs there are such that $i+j \leq n$. Set $i$ to $0$: $j$ can range from $0$ to $n$, then set $i$ to $1$: $j$ can range from $0$ to $n-1$, etc. The total number of pairs is thus $n+1+n+n-1+\dots+1 = \frac{(n+1)(n+2)}{2}$. Therefore the left summation evaluates to 
 \[\frac{(2n+1)(n+1)(n+2)}{2}\]

 Now we calculate $\sum_{i+j \leq n}^{} i+2j$. This simplifies to $\sum_{i+j \leq n}^{} i + 2 \cdot \sum_{i+j \leq n}^{} j$. Note that because $i+j = n$ is symmetric with respect to $i,j$, the sum of $i$ in all of the pairs will be equal to the sum of $j$ in all of the pairs. Thus this is equal to calculating $3 \cdot \sum_{i+j \leq n}^{} i$.

 In the pairs, $i=1$ appears for $j$ ranging between $0$ and $n-1$ so the sum here is $1 \cdot (n)$. Similarly $i=2$ appears for $j$ ranging from $0$ to $n-2$, so the sum is $2 \cdot (n-1)$. If we continue the pattern, the sum overall is $(n)+2 \cdot (n-1) + 3 \cdot (n-2) + \dots + (n) \cdot 1$. We can rearrange this as $((n)+(n-1)+ \dots + 1) + ((n-1)+(n-2)+ \dots + 1)+ ((n-2)+(n-3)+ \dots + 1) + \dots + 1)$

 \[= \frac{(n)(n+1)}{2} + \frac{(n-1)(n)}{2}+ \dots + 1\]

 We can write this in easier terms as $\sum_{k=0}^{n} \frac{(k)(k+1)}{2} = \frac{1}{2} \cdot \sum_{k=0}^{n} k^2+k$
 \[=\frac{1}{2} \cdot( \sum_{k=0}^{n} k^2 + \sum_{k=0}^{n} k)\]
 \[= \frac{1}{2} \cdot ( \frac{(n)(n+1)(2n+1)}{6} + \frac{(n)(n+1)}{2})\]
 \[= \frac{1}{2} \cdot ( \frac{(n)(n+1)(2n+1)}{6} + \frac{3n(n+1)}{6}) = \frac{1}{2} \cdot \frac{n(n+1)(2n+4)}{6}\]
 \[= \frac{n(n+1)(n+2)}{6}\]

 We multiply this by $3$ to obtain that 
 \[\sum_{i+j \leq n}^{} i+2j = \frac{n(n+1)(n+2)}{2}\]

 Thus our final answer for the number of distinct triplets $(h,i,j)$ is:
 \[\sum_{i+j\leq n}^{} 2n+1-i-2j = \frac{(2n+1)(n+1)(n+2)}{2} - \frac{n(n+1)(n+2)}{2}\]
 \[= \frac{(n+1)(n+2)}{2} \cdot (2n+1-n) = \frac{(n+1)(n+2)}{2} \cdot (n+1)\]
 \[= \frac{(n+1)^2(n+2)}{2}\]

 Now most of the work is done. We set this equal to $936$ and prime factorize. $936 = 12 \cdot 78 = 2^3 \cdot 3^2 \cdot 13$, so $(n+1)^2(n+2) = 936 \cdot 2 = 2^4 \cdot 3^2 \cdot 13$. Clearly $13$ cannot be anything squared and $2^4 \cdot 3^2$ is a perfect square, so $n+2 = 13$ and $n = 11 = \boxed{A}$




\textbf{Solution 2}

The product can be written as 
 \begin{align*} 2^a 3^b 4^c 5^d 6^e & = 2^{a + 2c + e} 3^{b + e} 5^d . \end{align*}


 Therefore, we need to find the number of ordered tuples $\left( a + 2c + e, b+e, d \right)$ where $a$, $b$, $c$, $d$, $e$ are non-negative integers satisfying $a+b+c+d+e \leq n$.We denote this number as $f(n)$.

 Denote by $g \left( k \right)$ the number of ordered tuples $\left( a + 2c + e, b+e \right)$ where $\left( a, b, c, e \right) \in \Delta_k$ with $\Delta_k \triangleq \left\{ (a,b,c,e) \in \Bbb Z_+^4: a+b+c+e \leq k \right\}$.

 Thus,
 \begin{align*} f \left( n \right) & = \sum_{d = 0}^n g \left( n - d \right) \\ & = \sum_{k = 0}^n g \left( k \right)  . \end{align*}


 Next, we compute $g \left( k \right)$.

 Denote $i = b + e$. Thus, for each given $i$, the range of $a + 2c + e$ is from 0 to $2 k - i$.Thus, the number of $\left( a + 2c + e, b + e \right)$ is
 \begin{align*} g \left( k \right) & = \sum_{i=0}^k \left( 2 k - i + 1 \right) \\ & = \frac{1}{2} \left( k + 1 \right) \left( 3 k + 2 \right) . \end{align*}


 Therefore, 
 \begin{align*} f \left( n \right) & = \sum_{k = 0}^n g \left( k \right)  \\ & = \sum_{k=0}^n \frac{1}{2} \left( k + 1 \right) \left( 3 k + 2 \right) \\ & = \frac{3}{2} \sum_{k=0}^n \left( k + 1 \right)^2 - \frac{1}{2} \sum_{k=0}^n \left( k + 1 \right) \\ & = \frac{3}{2} \cdot \frac{1}{6} \left( n+1 \right) \left( n+2 \right) \left( 2n + 3 \right) - \frac{1}{2} \cdot \frac{1}{2} \left( n + 1 \right) \left( n + 2 \right) \\ & = \frac{1}{2} \left( n + 1 \right)^2 \left( n + 2 \right) . \end{align*}


 By solving $f \left( n \right) = 936$, we get $n = \boxed{\textbf{(A) 11}}$.

 



\textbf{Solution 3 (Cheese)}

The product can be written as
 \begin{align*} 2^x 3^y 5^z \end{align*}


 Letting $n=1$, we get $(x,y,z)=(0,0,0),(0,0,1),(0,1,0),(1,0,0),(1,1,0),(2,0,0)$, 6 possible values. But if the only restriction of the product if that $2x\le n,y\le n,z\le n$, we can get $(2+1)(1+1)(1+1)=12$ possible values. We calculate the ratio
 \[r = \frac{\text{possible values of real situation}}{\text{possible values of ideal situation}} = \frac{6}{12}=0.5.\]

 Letting $n=2$, we get $(x,y,z)=(0,0,0),(0,0,1),(0,0,2),(0,1,0),(0,1,1),(1,0,0),(1,0,1),(1,1,0),(1,1,1),(1,2,0),$
$(2,0,0),(2,0,1),(2,1,0),(2,2,0),(3,0,0),(3,1,0),(4,0,0)$, 17 possible values.The number of possibilities in the ideal situation is $5*3*3=45$, making $r = 17/45 \approx 0.378$.

 Now we can predict the trend of $r$: as $n$ increases, $r$ decreases.Letting$n=3$, you get possible values of ideal situation=$7*4*4=112$.$n=4$, the number=$9*5*5=225$.
$n=5$, the number=$11*6*6=396$.
$n=6$, the number=$13*7*7=637,637<936$ so 6 is not the answer.
$n=7$, the number=$15*8*8=960$.
$n=8$, the number=$17*9*9=1377$,but $1377*0.378$\approx$521$ still much smaller than 936.
$n=9$, the number=$19*10*10=1900$,but $1900*0.378$\approx$718$ still smaller than 936.
$n=10$, the number=$21*11*11=2541$, $2541*0.378$\approx$960$ is a little bigger 936, but the quotient of (possible values of real situation)/(possible values of ideal situation) is much smaller than 0.378 now, so 10 is probably not the answer,so the answer is $\boxed{\textbf{(A) 11}}$.

 Check calculation:$n=11$,the number=$23*12*12=3312$,$3312*0.378$\approx$1252$ is much bigger than 936.

 



\textbf{Solution 4 (Easy computation)}

The key observation is if $P=2^a\cdot 3^b \cdot 5^c$, then given $b$ and $c$, $a$ can take any value from $0$ to $b+2d$ where $d=n-b-c$ is the number of rolls which is neither divisible by $3$ nor $5$. We are left to calculate
 \[\sum_{b+c\leq n} (b+2d+1)=\sum_{b+c+a=n} (b+2d+1).\]By symmetry, $\sum_{b+c+d=n} d = \sum_{b+c+d=n} c$. Therefore,
 \[\sum_{b+c\leq n}(b+2d+1)=\sum_{b+c\leq n}(b+c+d+1)=\sum_{b+c\leq n}(n+1)=(n+1)\binom{n+2}{2}.\]

 The rest is the same as above.

 



\textbf{Solution 5 ("normalize" the product) by FireSirius}

The difficulty is that a product can be formed by different combinations of 1,2,...,n. So the key thought of this solution is to set a "normalized" combination for each product.

 $P=1^{k_1}\cdot 2^{k_2} \cdot 3^{k_3}\cdot 4^{k_4}\cdot 5^{k_5} \cdot 6^{k_6}$, with $k_1+k_2+...+k_6=n$.

 For a fixed number of dice, it is easy to see that there are three basic "transformation": 1*6 = 2*3, 1*4 = 2*2, and 2*6 = 3*4, which does not change the number of dice. Other transformation like 1*6*6=3*3*4 can be viewed as a 1*6=2*3 combined with a 2*6=3*4. 

 So we can transform any given combination into a "normallised" combination in this way: step 1, if there are 1 and 6, then transorm them into 2*3; step 2, if there are both 1 and 4, then transform them into a 2*2; step3: if there are both 2 and 6, then transform them into 3*4. Until no more basic transformation can be made.

 Now begin casework with $k_1$:

 case 1($k_1 <= k_6$): all "1"s cancel with "6"s, so the product contains only 2,3,4,5,6. And now:

 If "2"s are no more than remaining "6"s: all "2"s cancel with "6"s, so the "nomarlized" product consists of only $"3456"$, with the total number of dice unchanged. No transformation can be made now and it is the same with n stars and 3 bars, so there are $C^3_{n+3}$ ways. 

 If "2"s are no less than remaining "6"s: all "6" cancels with "2", so the product consist of only $"2345"$. Again, there are $C^3_{n+3}$ ways.

 case 2($k_1>=k_6$): all "6" cancels with "1", the remaining numbers are: 1,2,3,4,5; now if "1"s are no less than "4"s, then it can be transformed into product with only $"1235"$; there are, again, $C^3_{n+3}$ ways. if "1"s are no more than "4", then it is product with only 2,3,4,5, which is already counted before.

 So there are three major cases:  $"2345"$, $"3456"$ and $"1235"$, each with $C^3_{n+3}$ ways. There are overlaps: $3^{k_3} 4^{k_4} 5^{k_5}$, which has $C^2_{n+2}$ ways, are counted in both $"2345"$ and $"3456"$, and $2^{k_2} 3^{k_3} 5^{k_5}$  are counted in both $"2345"$ and $"1235"$.

 So, the final number is: 
 \[3C^3_{n+3}-2C^2_{n+2}=\frac{(n+1)^2 (n+2)}{2},\]and now any possible product is counted exactly once. And then it is easily to find out $n = \boxed{\textbf{(A) 11}}$.

 



\textbf{Solution 6 (defining sequences)}

We define:

 $a_n$ is the number of possible products of the results of $n$ dice being rolled.

 
Now consider a modified five-sided die that excludes $5$, having faces $1$, $2$, $3$, $4$, and $6$. Let the number of possible products of the results of $n$ of these modified dice be equal to $b_n$.

 If in the $n$ rolls of the dice, we roll $5$ zero times, then the number of distinct products is equal to $b_n$. If we roll $5$ once, then the number of distinct products is equal to $b_{n-1}$. In general, if we roll $5$ $i$ times, then the number of distinct products is equal to $b_{n-i}$. Between the different cases of how many times we roll $5$, the products are guaranteed to be distinct. This is because in their prime factorizations, the powers of $5$ are different, and by the Fundmamental Theorem of Arithmetic, the products must be different as well.

 In other words, we have found:
 \[a_n=\sum_{i=0}^{n} b_i\]

 Now let us focus on $b_n$, and from now on, until otherwise specified, the die will refer to the modified, 5-sided die.In these $n$ rolls of the dice, if we roll $3$ zero times and $6$ zero times, then the product must be a power of $2$. It can range from $2^0$ to $4^n$, which is equal to $2^{2n}$. Therefore, there are $2n+1$ possible products. If we either roll $3$ once or $6$ once (not both), then the product must be $3$ times a power of $2$. This power of $2$ can range from $2^0$ to $2^{2(n-1)+1}=2^{2n-1}$, when we roll one $6$ and all other rolls are $4$. Therefore, there are $2n$ products. In general, if we roll $3$ and $6$ a combined $i$ times, then the product must be in the form $2^k * 3^i$, and $k$ ranges from $0$ to $2(n-i)+i=2n-i$, resulting in $2n-i+1$ products.

 In other words, we have found:
 \[b_n=\sum_{i=0}^{n} 2n-i+1\]

 We can now calculate $b_n$. Writing out the full sum, we get:
 \[b_n=(2n+1)+(2n)+\cdots+(n+1)\]This simplifies to:
 \[b_n=\frac{(n+1)(3n+2)}{2}=\frac{3n^2+5n+2}{2}\]

 Now, finally, we can calculate $a_n$:
 \[a_n=\sum_{i=0}^{n} b_i\]

 \[=\sum_{i=0}^{n} \frac{3i^2+5i+2}{2}\]

 \[=\sum_{i=0}^{n} \frac{3i^2}{2} + \sum_{i=0}^{n} \frac{5i}{2} + \sum_{i=0}^{n} 1\]

 \[=\frac{3}{2}*\frac{n(n+1)(2n+1)}{6} + \frac{5}{2}*\frac{n(n+1)}{2} + (n+1)\]

 \[=(n+1)(\frac{n(2n+1)}{4}+\frac{5n}{4}+1)\]

 \[=(n+1)(\frac{2n^2+6n+4}{4})\]

 \[=\frac{(n+1)^2(n+2)}{2}\]

 Finally, plugging in $a_n=936$, we get that $n=\boxed{\textbf{(A) 11}}$.

 


	\item (\textit{AMC 12 12B 2023}) Suppose that $a$, $b$, $c$ and $d$ are positive integers satisfying all of the following relations.


 
 \[abcd=2^6\cdot 3^9\cdot 5^7\]

 \[\text{lcm}(a,b)=2^3\cdot 3^2\cdot 5^3\]

 \[\text{lcm}(a,c)=2^3\cdot 3^3\cdot 5^3\]

 \[\text{lcm}(a,d)=2^3\cdot 3^3\cdot 5^3\]

 \[\text{lcm}(b,c)=2^1\cdot 3^3\cdot 5^2\]

 \[\text{lcm}(b,d)=2^2\cdot 3^3\cdot 5^2\]

 \[\text{lcm}(c,d)=2^2\cdot 3^3\cdot 5^2\]


 What is $\text{gcd}(a,b,c,d)$?


 $\textbf{(A)}~30\qquad\textbf{(B)}~45\qquad\textbf{(C)}~3\qquad\textbf{(D)}~15\qquad\textbf{(E)}~6$

 \textbf{Answer: }C



\textbf{Solution 1}

Denote by $\nu_p (x)$ the number of prime factor $p$ in number $x$.

 We index Equations given in this problem from (1) to (7).

 
First, we compute $\nu_2 (x)$ for $x \in \left\{ a, b, c, d \right\}$.

 Equation (5) implies $\max \left\{ \nu_2 (b), \nu_2 (c) \right\} = 1$.Equation (2) implies $\max \left\{ \nu_2 (a), \nu_2 (b) \right\} = 3$.Equation (6) implies $\max \left\{ \nu_2 (b), \nu_2 (d) \right\} = 2$.Equation (1) implies $\nu_2 (a) + \nu_2 (b) + \nu_2 (c) + \nu_2 (d) = 6$.

 Therefore, all above jointly imply $\nu_2 (a) = 3$, $\nu_2 (d) = 2$, and $\left( \nu_2 (b), \nu_2 (c) \right) = \left( 0 , 1 \right)$ or $\left( 1, 0 \right)$.

 
Second, we compute $\nu_3 (x)$ for $x \in \left\{ a, b, c, d \right\}$.

 Equation (2) implies $\max \left\{ \nu_3 (a), \nu_3 (b) \right\} = 2$.Equation (3) implies $\max \left\{ \nu_3 (a), \nu_3 (c) \right\} = 3$.Equation (4) implies $\max \left\{ \nu_3 (a), \nu_3 (d) \right\} = 3$.Equation (1) implies $\nu_3 (a) + \nu_3 (b) + \nu_3 (c) + \nu_3 (d) = 9$.

 Therefore, all above jointly imply $\nu_3 (c) = 3$, $\nu_3 (d) = 3$, and $\left( \nu_3 (a), \nu_3 (b) \right) = \left( 1 , 2 \right)$ or $\left( 2, 1 \right)$.

 Third, we compute $\nu_5 (x)$ for $x \in \left\{ a, b, c, d \right\}$.

 
Equation (5) implies $\max \left\{ \nu_5 (b), \nu_5 (c) \right\} = 2$.Equation (2) implies $\max \left\{ \nu_5 (a), \nu_5 (b) \right\} = 3$.Thus, $\nu_5 (a) = 3$.

 From Equations (5)-(7), we have either $\nu_5 (b) \leq 1$ and $\nu_5 (c) = \nu_5 (d) = 2$, or $\nu_5 (b) = 2$ and $\max \left\{ \nu_5 (c), \nu_5 (d) \right\} = 2$.

 Equation (1) implies $\nu_5 (a) + \nu_5 (b) + \nu_5 (c) + \nu_5 (d) = 7$.Thus, for $\nu_5 (b)$, $\nu_5 (c)$, $\nu_5 (d)$, there must be two 2s and one 0.

 Therefore,
 \begin{align*} {\textrm gcd} (a,b,c,d) & = \Pi_{p \in \{ 2, 3, 5\}} p^{\min\{ \nu_p (a), \nu_p(b) , \nu_p (c), \nu_p(d) \}} \\ & = 2^0 \cdot 3^1 \cdot 5^0 \\ & = \boxed{\textbf{(C) 3}}. \end{align*}





\textbf{Solution 2 (GCD/LCM Comparison)}

We are given that $abcd = 2^6 \cdot 3^9 \cdot 5^7$, and several LCM relations involving pairs of these numbers. Notice that

 $\mathrm{lcm}(a,b) \cdot \mathrm{lcm}(c,d) = (2^3 \cdot 3^2 \cdot 5^3) \cdot (2^2 \cdot 3^3 \cdot 5^2) = 2^{5} \cdot 3^{5} \cdot 5^{5}.$

 Comparing this product with $abcd$, we see that

 $abcd = 2^6 \cdot 3^9 \cdot 5^7 = \mathrm{lcm}(a,b) \cdot \mathrm{lcm}(c,d) \cdot 2^1 \cdot 3^4 \cdot 5^2.$

 This additional factor $2^1 \cdot 3^4 \cdot 5^2$ must be accounted for by the overlaps among $a$, $b$, $c$, $d$, which are their common divisors.

 The key observation is that the only prime with a leftover exponent greater than 3 is 3 (specifically, $3^4$), which suggests that all four numbers share at least one factor of 3.

 For primes 2 and 5, the leftover exponents are too small (1 and 2, respectively) to be shared by all four numbers, since the GCD exponent must be the minimum exponent among $a$, $b$, $c$, $d$.

 Thus, the only possible nontrivial common factor among $a$, $b$, $c$, $d$ is 3.

 Therefore,

 
$\gcd(a,b,c,d) = \boxed {3}.$



	\item (\textit{AMC 12 12B 2023}) A regular pentagon with area $\sqrt{5}+1$ is printed on paper and cut out. The five vertices of the pentagon are folded into the center of the pentagon, creating a smaller pentagon. What is the area of the new pentagon?


 $\textbf{(A)}~4-\sqrt{5}\qquad\textbf{(B)}~\sqrt{5}-1\qquad\textbf{(C)}~8-3\sqrt{5}\qquad\textbf{(D)}~\frac{\sqrt{5}+1}{2}\qquad\textbf{(E)}~\frac{2+\sqrt{5}}{3}$

 \textbf{Answer: }B



\textbf{Solution 1}

Since $A$ is folded onto $O$, $AM = MO$ where $M$ is the intersection of $AO$ and the creaseline between $A$ and $O$. Note that the inner pentagon is regular, and therefore similar to the original pentagon, due to symmetry.

 Because of their similarity, the ratio of the inner pentagon's area to that of the outer pentagon can be represented by 

 $\left(\frac{OM}{ON}\right)^{2} = \left(\frac{\frac{OA}{2}}{OA\sin (\angle OAE)}\right)^{2} = \frac{1}{4\sin^{2}54}$ 

 


 \textbf{Option 1: Knowledge} Remember that $\sin54 = \frac{1+\sqrt5}{4}$.

 \textbf{Option 2: Angle Identities} $\cos54 = \sin36$
$4\cos^{3}18-3\cos18 = 2\sin18\cos18$
$4(1-\sin^{2}18)-3-2\sin18=0$
$4\sin^{2}18+2\sin18-1=0$
$\sin18 = \frac{-1+\sqrt5}{4}$
$\sin54 = \cos36 = 1-2\sin^{2}18 = \frac{1+\sqrt5}{4}$

 \textbf{Option 3: Golden Ratio/Triangle} It is well known that a 36-36-108 triangle has side lengths with ratio 1 : 1 : $\varphi$

 Dropping the altitude to the longer side results in $\sin 54 = \frac{\varphi}{2} = \frac{\sqrt{5} + 1}{4}$ (\href{https://artofproblemsolving.com/wiki/index.php?title=User:Elib0128&action=edit&redlink=1}{Elib0128} (\href{https://artofproblemsolving.com/wiki/index.php?title=User\_talk:Elib0128&action=edit&redlink=1}{talk}))

 
$\sin^{2}54 =\frac{3+\sqrt5}{8}$

 Let the inner pentagon be $Z$.

 $[Z] = \frac{1}{4\sin^{2}54}[ABCDE]$
$= \frac{2(1+\sqrt5)}{3+\sqrt5}$
$= \sqrt5-1$

 So the answer is $\boxed{B}$

 -Dissmo Thegoat

 



\textbf{Solution 2}


                    \begin{figure}[H]
                        \centering
                        \includegraphics[width=0.5\textwidth]{Solution/AMC_12_12B_2023_25_Solution_Figure_1.png}
                    \end{figure}
                    

 We can find the area of the red pentagon by taking the area of the total pentagon and subtracting the area outside the red pentagon. 

 The area outside the red pentagon is the sum of the larger isosceles triangles, but this double counts the overlapping regions of the small isosceles triangles, so we have to subtract those out.

 We have $[FGHIJ] = [ABCDE]-(5 \cdot[DKT]-5 \cdot [PFK])$

 Lets focus on finding the area of each individual triangle:

 
                    \begin{figure}[H]
                        \centering
                        \includegraphics[width=0.5\textwidth]{Solution/AMC_12_12B_2023_25_Solution_Figure_2.png}
                    \end{figure}
                    

 Notice that we have no information about the side length, so instead we let the side length be $s$. Now we can drop an altitude from $O$ to the base of the triangle, and we know this altitude must split the base of the pentagon in half, so we can create a right triangle. Furthermore, draw a line from $O$ to $D$. This must bisect angle $D$ which is $108$ degrees, so we create $36-54-90$ triangles. Specifically, we know $\angle ODK = 54^{\circ}$, $\angle DOU = 36^{\circ}$, and $\angle DTK = 36^{\circ}$ because $\triangle DTK$ is isosceles and we know the vertex angle is $108^{\circ}$. We encode this information in the diagram below:
                    \begin{figure}[H]
                        \centering
                        \includegraphics[width=0.5\textwidth]{Solution/AMC_12_12B_2023_25_Solution_Figure_3.png}
                    \end{figure}
                    

 Since $\triangle DKT$ is isosceles, the area of $\triangle DVT$ is half the area of $\triangle DKT$. Similarly, the area of $\triangle UFK$ is half that of $\triangle PFK$. Thus:

 

 \[[FGHIJ] = [ABCDE]-(5 \cdot[DKT]-5 \cdot [PFK]) \implies [FGHIJ] = [ABCDE]-(10 \cdot [DVT]-10 \cdot [UFK])\]

 We also know that since we dropped an altitude from $O$ to $U$, the area of $\triangle ODU$ must be half of a fifth of the total area of the pentagon. Therefore we can rewrite the above equation as 
 \[[FGHIJ]=10 \cdot ([ODU]-[DVT]+[UFK])\]

 Now notice that $\triangle ODU  ~ \triangle TDV ~ \triangle KFU$ by AA similarity. Therefore, if we can write the areas of the latter two triangles as a ratio of the first triangle, we can express the whole equation in terms of $[ODU]$, and by extension $[ABCDE]$, which we know. To find these ratios, we can find the side length ratios and square them because the triangles are similar.

 We already know $DU = \frac{s}{2}$, so let's try to find it's analogous side for $\triangle TDV$ and $\triangle KFU$. These sides are $DV$ and $FU$, respectively. 

 First, $\frac{s}{2} = OD \cdot cos(54^{\circ})$, so $OD = \frac{s}{2} \cdot sec(54^{\circ})$. Then notice that $DV = \frac{OD}{2}$ because we have to fold $D$ to hit $O$, so the folding crease has to be exactly halfway between $O$ and $D$. Therefore, 
 \[DV = \frac{s}{4} \cdot sec(54^{\circ}) \implies \frac{DU}{DV} = \frac{\frac{s}{2}}{\frac{s}{4} \cdot sec(54^{\circ})} = 2 \cdot cos(54^{\circ})\]

 Now the ratio between the area of two similar triangles is the square of the ratio of their analogous side lengths. Thus
 \[\frac{[ODU]}{[TDV]} = 4 \cdot cos^2(54^{\circ}) \implies \frac{10 \cdot [ODU]}{10 \cdot [TDV]} = 4 \cdot cos^2(54^{\circ})\]
 \[\implies \frac{[ABCDE]}{10 \cdot [TDV]} = 4 \cdot cos^2(54^{\circ}) \implies \frac{\sqrt{5}+1}{10 \cdot [TDV]} = 4 \cdot cos^2(54^{\circ})\]
 \[\implies 10 \cdot [TDV] = \frac{\sqrt{5}+1}{4 \cdot cos^2(54^{\circ})} \implies 10 \cdot [TDV] = \frac{\sqrt{5}+1}{4} \cdot sec^2(54^{\circ})\]

 Now let's move on and calculate the ratio of the other side length. Calculating $FU$ is slightly tricker. First, we find $TD$: $TD \cdot cos(54^{\circ}) = DV = \frac{s}{4} sec^54^{\circ}) \implies TD = \frac{s}{4} \cdot sec^2(54^{\circ})$. Now since $\triangle DTK$ is isosceles, $TD = DK$ and $UK = DK-DU = TD-DU = \frac{s}{4} \cdot sec^2(54^{\circ})-\frac{s}{2} = \frac{s}{4} \cdot (sec^2(54^{\circ})-2) = \frac{s}{4} \cdot (tan^2(54^{\circ})-1)$.

 Now $FU = UK \cdot tan(36^{\circ}) \implies FU \frac{s}{4} \cdot (tan^2(54^{\circ})-1) \cdot tan(36^{\circ})$. Now note that $tan(x) \cdot tan(90-x) = 1$ because opposite over adjacent cancel each other out in a right triangle. Thus, $FU = \frac{s}{4} \cdot (tan(54^{\circ})-tan(36^{\circ}))$ 

 Now, 
 \[\frac{DU}{FU} = \frac{\frac{s}{2}}{\frac{s}{4} \cdot (tan(54^{\circ})-tan(36^{\circ}))} = \frac{2}{tan(54^{\circ})-tan(36^{\circ})}\]

 \[\implies \frac{[DU]}{[FU]} = \frac{4}{(tan(54^{\circ})-tan(36^{\circ}))^2} \implies 10 \cdot [FU] = \frac{1+\sqrt{5}}{4} \cdot (tan(54^{\circ})-tan(36^{\circ}))^2\]

 Now we go back to our first equation and plug in our values:

 
 \[[FGHIJ] = [ABCDE]-(10 \cdot [DVT]-10 \cdot [UFK]) \implies [FGHIJ] = 1+\sqrt{5}-\frac{\sqrt{5}+1}{4} \cdot sec^2(54^{\circ})+\frac{1+\sqrt{5}}{4} \cdot (tan(54^{\circ})-tan(36^{\circ}))^2\]

 \[\implies [FGHIJ] = \frac{1+\sqrt{5}}{4} \cdot (4-sec^2(54^{\circ})+(tan(54^{\circ})-tan(36^{\circ}))^2)\]

 Note $(tan(x)-tan(90-x))^2 = tan^2(x)-2 \cdot tan(x) \cdot tan(90-x)+tan^2(90-x) = tan^2(x)-2+tan^2(90-x)$.

 Also note that $tan^2(x)+1 = sec^2(x)$.Thus 
 \[[FGHIJ] = \frac{1+\sqrt{5}}{4} \cdot (2-sec^2(54^{\circ})+tan^2(54^{\circ})+tan^2(36^{\circ}))\]
 \[\implies [FGHIJ] = \frac{1+\sqrt{5}}{4} \cdot (2+(-1)+tan^2(36^{\circ})) \implies [FGHIJ] = \frac{1+\sqrt{5}}{4} \cdot (1+tan^2(36^{\circ}))\].

 Now all that remains is to find $tan^2(36^{\circ})$. We can use the tan addition formula to find the general form of $tan(5x)$ or remember question 25 from this year's AMC 12A. We have that 
 \[tan(5x) = \frac{5tan(x)-10tan^3(x)+tan^5(x)}{1-10tan^2(x)+5tan^4(x)}\]. 

 Plug in $x=36$. Then we have 
 \[tan(180^{\circ}) = \frac{5tan(36^{\circ})-10tan^3(36^{\circ})+tan^5(36^{\circ})}{1-10tan^2(36^{\circ})+5tan^4(36^{\circ})}\]Now let $y = tan(36^{\circ})$. We have the equation 
 \[\frac{5y-10y^3+y^5}{1-10y^2+5y^4} = 0 \implies 5y-10y^3+y^5 = 0\]
 \[\implies 5-10y^2+y^4 = 0 \implies 5-10z+z^2 = 0\]Where we let $z = y^2$. Using the quadratic formula, we have 
 \[z = \frac{10 \pm \sqrt{80}}{2} = 5 \pm 2\sqrt{5}\]Now since $y = tan(36^{\circ})$, $z = tan^2(36^{\circ})$, which is what we were looking for. Notice that $tan(0^{\circ}) = 0$ and $tan(45^{\circ}) = 1$, so $tan(36^{\circ})$ is between $0$ and $1$, and so is it's square. Thus $z = 5 - 2\sqrt{5}$, not the other root.

 Finally:

 
 \[[FGHIJ] = \frac{1+\sqrt{5}}{4} \cdot (1+tan^2(36^{\circ})) = \frac{1+\sqrt{5}}{4} \cdot (1+5-2\sqrt{5})\]
 \[\implies [FGHIJ] = \frac{1+\sqrt{5}}{4} \cdot (6-2\sqrt{5}) = \frac{4\sqrt{5}-4}{4} = \sqrt{5}-1\]

 Therefore, 
 \[[FGHIJ] = \sqrt{5}-1 = \boxed{B}\]




\textbf{Solution 3}


                    \begin{figure}[H]
                        \centering
                        \includegraphics[width=0.5\textwidth]{Solution/AMC_12_12B_2023_25_Solution_Figure_4.png}
                    \end{figure}
                    

 Let $r_b$ and $r_s$ be the circumradius of the big and small pentagon, respectively. Let $a_s$ be the apothem of the smaller pentagon and $A_s$ and $A_b$ be the areas of the smaller and larger pentagon, respectively. 

 From the diagram:
 \begin{align*}     \cos{36^\circ} &= \dfrac{a_s}{r_s} = \dfrac{\phi}{2} = \dfrac{\sqrt{5}+1}{4}\\     a_s &= \dfrac{r_b}{2}\\     A_s &= \left(\dfrac{r_s}{r_b}\right)^2A_b\\     &=\left(\dfrac{a_s}{\cos{36^\circ} r_b}\right)^2 \left(1+\sqrt{5}\right)\\     &=\left(\dfrac{a_s}{\dfrac{\phi}{2} r_b}\right)^2 \left(1+\sqrt{5}\right)\\     &=\left(\dfrac{1}{2 \dfrac{\phi}{2}}\right)^2 \left(1+\sqrt{5}\right)\\     &=\left(\dfrac{2}{\sqrt{5}+1}\right)^2 \left(1+\sqrt{5}\right)\\     &=\dfrac{4}{\sqrt{5}+1} \\     &=\dfrac{4\left(\sqrt{5}-1\right)}{\left(\sqrt{5}+1\right)\left(\sqrt{5}-1\right)} \\     &=\sqrt{5}-1 \end{align*}

 \[\boxed{\textbf{(B) }\sqrt{5}-1}\]~Technodoggo

 



\textbf{Solution 4}

Interestingly, we find that the pentagon we need is the one that is represented by the intersection of perpendicular bisectors of the connection from the center of the pentagon to one vertex. Through similar triangles and the golden ratio, we find that the side length ratio of the two pentagons is $\frac{\sqrt{5}-1}{2}$ Thus, the answer is $(\sqrt{5}+1) \cdot (\frac{\sqrt{5}-1}{2})^2 = \sqrt{5}-1$. $\boxed{\text{B}}$~andliu766

 \textbf{Supplement (Calculating sin54/cos36 from Scratch)} Method 1:

 
                    \begin{figure}[H]
                        \centering
                        \includegraphics[width=0.5\textwidth]{Solution/AMC_12_12B_2023_25_Solution_Figure_5.png}
                    \end{figure}
                    Construct golden ratio triangle $\triangle ABC$ with $\angle A = 36^{\circ}$, $\angle B = \angle C = 72^{\circ}$ and $\triangle BCD$ with $\angle C = 36^{\circ}$, $\angle DBC = \angle BDC = 72^{\circ}$. WLOG, let $AB = AC = 1$, $BC = CD = AD = a$, $BD = 1-a$. $\triangle BAC \sim \triangle BCD$

 \[\frac{AC}{BC} = \frac{BC}{BD}, \quad \frac{1}{a} = \frac{a}{1-a}, \quad 1-a=a^2, \quad a^2 + a - 1 = 0\]

 \[a = \frac{ -1 + \sqrt{1^2  - 4(-1) } }{2} = \frac{ \sqrt{5} -1 }{2}\]

 \[\cos 36^{\circ} = \cos \angle A = \frac{AE}{AC} = \frac{ 1-a }{2} + a = \frac{ a + 1 }{2} = \frac{ \frac{ \sqrt{5} -1 }{2} + 1 }{2} = \frac{ \sqrt{5} + 1 }{4}\]

 \[\sin 54^{\circ} = \cos 36^{\circ} = \frac{ \sqrt{5} + 1 }{4}\]

 Method 2:

 As explained \href{https://artofproblemsolving.com/wiki/index.php/2022\_AMC\_12B\_Problems/Problem\_24\#Supplement\_.28Explanation\_of\_why\_cos.282.CF.80.2F7.29\_.2B\_cos.284.CF.80.2F7.29\_.2B\_cos.286.CF.80.2F7.29\_.3D\_-1.2F2.29}{here}, $\cos \frac{2 \pi}{5} + \cos \frac{4 \pi}{5} = - \frac12$

 \[\cos \frac{2 \pi}{5} - \cos \frac{\pi}{5} = - \frac12\]

 \[2(\cos\frac{ \pi}{5})^2 - 1 - \cos \frac{\pi}{5} = -1/2\]

 \[4(\cos \frac{\pi}{5})^2 - 2 \cos \frac{\pi}{5} - 1 = 0\]

 \[\cos 36^{\circ} = \cos \frac{\pi}{5} = \frac{2 + \sqrt{2^2 + 4 \cdot 4} }{8} = \frac{1+\sqrt{5}}{4}\]




\textbf{Solution 5 (Much easier if you notice the one line construction + use properties of the golden ratio)}

We will start with a similar diagram to all of the other submitted solutions with ONE difference: We will connect line segment $CR$ to form a $36^\circ$-$36^\circ$-$108^\circ$ golden triangle $\triangle CPR$. We know this is a $36^\circ$-$36^\circ$-$108^\circ$ golden triangle because $QR$ is the perpendicular bisector of $CP$ (in order to fold vertex $C$ onto the center $P$), so $\triangle CPR$ must be isosceles, with angle $\angle CPR = \frac{360^\circ}{10} = 36^\circ$ due to the regular pentagon.

 
                    \begin{figure}[H]
                        \centering
                        \includegraphics[width=0.5\textwidth]{Solution/AMC_12_12B_2023_25_Solution_Figure_6.png}
                    \end{figure}
                    

                    \begin{figure}[H]
                        \centering
                        \includegraphics[width=0.5\textwidth]{Solution/AMC_12_12B_2023_25_Solution_Figure_7.png}
                    \end{figure}
                    

 Rather than trigonometry, it is well known that $\frac{CP}{PR} = \frac{\phi}{1}$ due to the golden triangle. Since both pentagons are regular, they are similar. Therefore the ratio of the areas is equal to the ratio of the radii squared (rather than using the ratio of the sides of the pentagon). Let $x$ be the area of the smaller pentagon. Note that we are given the larger pentagon has area $2\phi$. 

 We will also use a well known identity of the golden ratio: $\frac{1}{\phi} = \phi - 1$ (and for completeness sake for any student reading this in the future looking for other golden ratio identities: $\phi^2 = \phi + 1$) to simplify expressions involving the golden ratio rather than to do the algebra on the square roots. Then:

 
 \begin{align*} \frac{x}{2\phi} &= \left(\frac{PR}{CP}\right)^2 \\ \frac{x}{2\phi} &= \left(\frac{1}{\phi}\right)^2 \\ \frac{x}{2} &= \frac{1}{\phi} \\ \frac{x}{2} &= \phi - 1 \\ x &= 2\phi - 2 \\ x &= 2\left(\frac{1+\sqrt{5}}{2}\right)-2 \\ x &= 1+\sqrt{5}-2 \\ x &= \sqrt{5}-1 \end{align*}

$\boxed{\text{B}}$



\end{enumerate}